\documentclass[UTF8, 12pt, fontset=fandol, oneside]{ctexart}
\usepackage{researchnotes}

\graphicspath{{math/概率论/}}

\title{第 2 章\quad 条件概率}
\author{}
\date{}

\begin{document}

\maketitle

我们已经把概率介绍为一种语言，用来表达我们对事件的信念程度，或表达我们面对事件时的不确定性。每当我们观察到新的证据（也就是获得数据）时，我们就获得了可能影响这种不确定性的信息。一个与既有信念相符的新观察，可能使我们更加确信这个信念；而一个出乎意料的观察，则可能使这个信念受到质疑。\textbf{条件概率}所要处理的，正是这样一个基本问题：在观察到证据之后，我们应该怎样更新自己的信念？

\section*{2.1\quad 条件思考的重要性}

条件概率对于科学、医学和法律推理都是基本的，因为它说明了如何以逻辑一致的方式把证据纳入我们对世界的理解。事实上，一个有用的观点是：\textit{所有概率都是条件的}；无论是否明确写出，每个概率背后都包含背景知识或假设。

例如，设某天早晨我们关心事件 $R$：当天会下雨。令 $P(R)$ 表示在向窗外看之前，我们对下雨概率的评估。如果随后向外看，发现天空乌云密布，那么下雨的概率大概应当增加；这个新的概率记作 $P(R|C)$（读作“给定 $C$ 时 $R$ 的概率”），其中 $C$ 表示天空乌云密布这一事件。当我们从 $P(R)$ 变为 $P(R|C)$ 时，称为“以 $C$ 为条件”。随着一天推进，我们还可能获得越来越多关于天气状况的信息；如果观察到事件 $B_1,\ldots,B_n$ 已经发生，那么在这些证据下下雨的条件概率可写作 $P(R|B_1,\ldots,B_n)$。如果最终观察到真的开始下雨，那么这个条件概率就变为 $1$。

此外，我们将会看到，条件化也是一种非常强有力的解题策略：它常常使我们能够把一个复杂问题分解成若干较易处理的情形，然后逐情形分析。正如在计算机科学中，常见策略是把一个大问题拆成较小的片段，在概率中，常见策略则是把复杂的概率问题化为一组较简单的条件概率问题。特别地，我们会讨论一种称为\textbf{第一步分析}的策略；对于实验具有多个阶段的问题，它常常可以给出递归解。

由于条件化既是根据证据更新信念的方法，又是重要的解题策略，我们说：
\[
  \text{\textit{条件化是统计学的灵魂。}}
\]

\section*{2.2\quad 定义与直觉}

\noindent\textbf{定义 2.2.1（条件概率）}\quad 若 $A$ 与 $B$ 是事件，且 $P(B)>0$，则给定 $B$ 时 $A$ 的\textbf{条件概率}记作 $P(A|B)$，定义为
\[
  P(A|B)=\frac{P(A\cap B)}{P(B)}.
\]

这里，$A$ 是我们想要更新其不确定性的事件，而 $B$ 是我们观察到的证据（或想要作为条件处理的信息）。我们把 $P(A)$ 称为 $A$ 的\textbf{先验概率}，把 $P(A|B)$ 称为 $A$ 的\textbf{后验概率}；“先验”表示在依据证据更新之前，“后验”表示在依据证据更新之后。

理解条件竖线之后的事件非常重要：它表示我们已经观察到、或正在作为条件的证据。$P(A|B)$ 是给定证据 $B$ 时 $A$ 的概率，而不是某个叫作 $A|B$ 的对象的概率。正如 2.4.1 节中的生物危害提示会讨论的，并不存在 $A|B$ 这样一个事件。

对任意事件 $A$，
\[
  P(A|A)=\frac{P(A\cap A)}{P(A)}=1.
\]
在观察到 $A$ 已经发生之后，我们对 $A$ 的更新概率就是 $1$。如果不是这样，我们就应当要求重新定义条件概率了。

\noindent\textbf{例 2.2.2（两张牌）}\quad 将一副标准扑克牌充分洗匀后，不放回地随机抽取两张牌，每次抽一张。令 $A$ 为第一张牌是红桃的事件，令 $B$ 为第二张牌是红色的事件。求 $P(A|B)$ 与 $P(B|A)$。

\noindent\textbf{解}\quad 由朴素概率定义与乘法规则，
\[
  P(A\cap B)=\frac{13\cdot 25}{52\cdot 51}=\frac{25}{204},
\]
因为有利结果可由任取 $13$ 张红桃之一作为第一张牌，再从剩余牌中任取 $25$ 张红色牌之一作为第二张牌确定。同时，$P(A)=1/4$，因为四种花色等可能；并且
\[
  P(B)=\frac{26\cdot 51}{52\cdot 51}=\frac12,
\]
因为第二张牌为红色有 $26$ 种可能，而对每一种第二张牌，第一张牌都可以是任意其他牌（回想第 1 章中的乘法规则：此处不需要按时间顺序计数）。

也可以用对称性更简洁地看出 $P(B)=1/2$：在实验进行之前，从整体角度看，第二张牌等可能地是牌堆中的任意一张牌。

现在所有使用条件概率定义所需的量都已经具备：
\[
  P(A|B)=\frac{P(A\cap B)}{P(B)}
  =\frac{25/204}{1/2}=\frac{25}{102},
  \qquad
  P(B|A)=\frac{P(B\cap A)}{P(A)}
  =\frac{25/204}{1/4}=\frac{25}{51}.
\]

这个例子很简单，但已经有几件事值得注意。

\begin{enumerate}
  \item 必须非常小心地区分哪些事件放在条件竖线的哪一侧。特别地，$P(A|B)\ne P(B|A)$。下一节会研究一般情形下 $P(A|B)$ 与 $P(B|A)$ 的关系。混淆这两个量称为\textbf{检察官谬误}，将在 2.8 节讨论。如果我们改为把 $B$ 定义成第二张牌是红桃，那么这两个条件概率就会相等。
  \item $P(A|B)$ 与 $P(B|A)$ 在直觉和数学上都有意义；抽牌的时间顺序并不决定我们可以考察哪一个条件概率。计算条件概率时，我们关心的是观察到一个事件能为另一个事件提供什么\textit{信息}，而不是一个事件是否\textit{导致}另一个事件。为了进一步获得直觉，可以想象有人把牌摊开，同时用左手抽一张牌、右手抽一张牌。若根据左手和右手的牌来定义 $A,B$，而不是根据第一张牌和第二张牌来定义，问题的结构并不会发生本质变化。
  \item 也可以直接解释 $P(B|A)=25/51$：如果第一张抽到的是红桃，那么剩余牌由 $25$ 张红牌与 $26$ 张黑牌组成（且下一张等可能地是其中任意一张），所以抽到红牌的条件概率是 $25/(25+26)=25/51$。用这种方式求 $P(A|B)$ 较困难：若知道第二张牌是红色，我们可能会想“这当然有用，但真正想知道的是它是不是红桃”。本章后面的条件概率结果会提供绕开这一困难的方法。
\end{enumerate}

为了进一步说明条件概率的含义，下面给出两个直观解释。

\noindent\textbf{直觉 2.2.3（鹅卵石世界）}\quad 考虑一个有限样本空间，把结果想象成总质量为 $1$ 的鹅卵石。由于 $A$ 是事件，它是一组鹅卵石；$B$ 也是如此。图 2.1(a) 给出了一个例子。

\begin{figure}[htbp]
  \centering
  \includegraphics[width=0.72\textwidth]{2-1.png}
  \caption{关于 $P(A|B)$ 的鹅卵石世界直觉。从左到右：(a) 事件 $A$ 与 $B$ 是样本空间的子集；(b) 由于我们知道 $B$ 已经发生，去掉 $B^c$ 中的结果；(c) 在受限样本空间中重新归一化，使总质量仍为 $1$。}
\end{figure}

现在假设我们得知 $B$ 已经发生。在图 2.1(b) 中，获得这一信息后，我们去掉所有 $B^c$ 中的鹅卵石，因为它们与 $B$ 已发生这一知识不相容。此时，$P(A\cap B)$ 是仍保留在 $A$ 中的鹅卵石总质量。最后，在图 2.1(c) 中，我们重新归一化，也就是把所有质量都除以同一个常数，使剩余鹅卵石的新总质量为 $1$。这个常数就是 $P(B)$，即 $B$ 中鹅卵石的总质量。对应于事件 $A$ 的更新后质量就是条件概率
\[
  P(A|B)=\frac{P(A\cap B)}{P(B)}.
\]

这样，我们的概率就按照观察到的证据得到了更新。与证据矛盾的结果被丢弃，其质量在剩余结果之间重新分配，同时保持剩余结果之间的相对质量不变。例如，若鹅卵石 2 初始质量是鹅卵石 1 的两倍，并且二者都包含在 $B$ 中，那么以 $B$ 为条件之后，鹅卵石 2 的质量仍是鹅卵石 1 的两倍。但若鹅卵石 2 不在 $B$ 中，则以 $B$ 为条件之后，它的质量更新为 $0$。

\noindent\textbf{直觉 2.2.4（频率学派解释）}\quad 回忆概率的频率学派解释：概率基于大量重复试验中的相对频率。设想把我们的实验重复许多次，得到一长串观察结果。给定 $B$ 时 $A$ 的条件概率可以自然地理解为：只关注 $B$ 发生的那些试验，在其中 $A$ 发生的比例。

在图 2.2 中，实验结果可以写成由 $0$ 和 $1$ 组成的字符串；$B$ 是第一位数字为 $1$ 的事件，$A$ 是第二位数字为 $1$ 的事件。以 $B$ 为条件时，我们把所有 $B$ 发生的重复试验圈出，然后观察在这些被圈出的重复试验中，事件 $A$ 也发生的比例。

用符号表示，令 $n_A,n_B,n_{AB}$ 分别为大量 $n$ 次重复实验中 $A,B,A\cap B$ 的发生次数。频率学派解释是
\[
  P(A)\approx \frac{n_A}{n},\qquad
  P(B)\approx \frac{n_B}{n},\qquad
  P(A\cap B)\approx \frac{n_{AB}}{n}.
\]
于是 $P(A|B)$ 可解释为 $n_{AB}/n_B$，它等于 $(n_{AB}/n)/(n_B/n)$。这个解释再次转化为
\[
  P(A|B)=\frac{P(A\cap B)}{P(B)}.
\]

\begin{figure}[htbp]
  \centering
  \includegraphics[width=0.58\textwidth]{2-2.png}
  \caption{关于 $P(A|B)$ 的频率学派直觉。$B$ 发生的重复试验被圈出；其中 $A$ 也发生的重复试验用粗体突出。$P(A|B)$ 是在 $B$ 发生的重复试验子集中，$A$ 发生的长期相对频率。}
\end{figure}

为了练习条件概率定义的应用，我们再看几个例子。接下来的三个例子都以一个有两个孩子的家庭为基本情形，但由于精确假设以及所条件化的信息不同，会出现一些细微差别。

\noindent\textbf{例 2.2.5（两个孩子）}\quad Martin Gardner 在 20 世纪 50 年代的《科学美国人》专栏中提出了下面的谜题。

\begin{quote}
Jones 先生有两个孩子。较年长的孩子是女孩。两个孩子都是女孩的概率是多少？

Smith 先生有两个孩子。其中至少一个是男孩。两个孩子都是男孩的概率是多少？
\end{quote}

乍看之下，这个问题似乎只是条件概率的简单应用，但几十年来，人们一直争论这两问是否、以及为什么应当有不同答案，也争论该问题在多大程度上是含混的。Gardner 给出的两个答案分别是 $1/2$ 和 $1/3$；这可能显得很悖论：为什么知道较年长孩子的性别，和只知道某个孩子的性别，会产生不同影响？

关键是必须澄清问题的假设。为了得到 Gardner 的答案，实际上使用了若干隐含假设。

\begin{itemize}
  \item 假设性别是二元的，因此每个孩子都能明确归为男孩或女孩。事实上，许多人并不能整齐地落入“男性”或“女性”这两个类别，并且会认同自己具有非二元性别。
  \item 假设对较年长孩子和较年幼孩子都有 $P(\text{男孩})=P(\text{女孩})$。事实上，在多数国家，男孩出生数略多于女孩。例如，在美国，通常估计每出生 $100$ 个女孩，就有约 $105$ 个男孩出生。
  \item 假设两个孩子的性别相互独立；也就是说，知道较年长孩子的性别，不会提供关于较年幼孩子性别的信息，反之亦然。如果两个孩子是同卵双胞胎之类的情形，这个假设就不现实。
\end{itemize}

在这些（诚然有问题的）简化假设下，问题可以如下求解。

\noindent\textbf{解}\quad 根据上面的假设，由条件概率定义，
\[
  P(\text{两个都是女孩}|\text{较年长的是女孩})
  =
  \frac{P(\text{两个都是女孩，较年长的是女孩})}
       {P(\text{较年长的是女孩})}
  =\frac{1/4}{1/2}=\frac12,
\]
\[
  P(\text{两个都是女孩}|\text{至少一个是女孩})
  =
  \frac{P(\text{两个都是女孩，至少一个是女孩})}
       {P(\text{至少一个是女孩})}
  =\frac{1/4}{3/4}=\frac13.
\]
为了便于比较两部分，我们把第二问也改用女孩来表述。两个结果不同可能显得违反直觉，因为似乎没有理由关心究竟是较年长孩子还是较年幼孩子为女孩。确实，由对称性，
\[
  P(\text{两个都是女孩}|\text{较年幼的是女孩})
  =
  P(\text{两个都是女孩}|\text{较年长的是女孩})
  =\frac12.
\]

然而，$P(\text{两个都是女孩}|\text{较年长的是女孩})$ 与 $P(\text{两个都是女孩}|\text{至少一个是女孩})$ 之间并没有这种对称性。说“较年长的孩子是女孩”指定了一个\textit{具体}孩子，然后另一个孩子（较年幼孩子）有 $50\%$ 的概率是女孩；“至少一个”并不指向某个具体孩子。以某个具体孩子是女孩为条件，会从样本空间 $\{GG,GB,BG,BB\}$ 中去掉 $2$ 个“鹅卵石”；这里例如 $GB$ 表示较年长孩子是女孩、较年幼孩子是男孩。相反，以至少一个孩子是女孩为条件，只会去掉 $BB$。

\noindent\textbf{例 2.2.6（随机遇到的孩子是女孩）}\quad 一个家庭有两个孩子。你随机遇到其中一个，并得知她是女孩。在与上一例相同的假设下，两个孩子都是女孩的条件概率是多少？再假设你等可能地遇到任一孩子，并且遇到哪一个孩子与性别无关。

\noindent\textbf{解}\quad 直观上，答案应当是 $1/2$：想象我们遇到的孩子就在面前，另一个孩子在家。两个孩子都是女孩，只是在说家里的那个孩子也是女孩；这似乎与面前这个孩子是女孩这一事实无关。不过我们仍用条件概率定义仔细检查，这也可以练习如何用集合记号写事件。

令 $G_1,G_2,G_3$ 分别表示较年长孩子、较年幼孩子、随机遇到的孩子是女孩的事件。由假设，
\[
  P(G_1)=P(G_2)=P(G_3)=\frac12.
\]
由朴素概率定义，或由 2.5 节中将解释的独立性，
\[
  P(G_1\cap G_2)=\frac14.
\]
因此
\[
  P(G_1\cap G_2|G_3)
  =
  \frac{P(G_1\cap G_2\cap G_3)}{P(G_3)}
  =
  \frac{1/4}{1/2}
  =
  \frac12,
\]
因为 $G_1\cap G_2\cap G_3=G_1\cap G_2$：若两个孩子都是女孩，那么随机遇到的孩子必然是女孩。

不过要记住，为了得到 $1/2$，我们需要关于随机孩子如何被选出的假设。用统计学语言说，我们收集了一个\textbf{随机样本}；这里样本由两个孩子中的一个组成。统计学中最重要的原则之一，就是必须仔细思考样本是如何收集的，而不能只盯着原始数据、却不了解它们来自哪里。举一个简单的极端情形，假设有一条压制性的法律规定：如果男孩有姐妹，他就不得离开家。那么“随机遇到的孩子是女孩”等价于“至少一个孩子是女孩”，于是问题就化为例 2.2.5 的第一部分。

\noindent\textbf{例 2.2.7（冬季出生的女孩）}\quad 一个家庭有两个孩子。已知两个孩子中至少有一个是冬季出生的女孩，求两个孩子都是女孩的概率。除例 2.2.5 的假设外，再假设四个季节等可能，并且性别与出生季节独立。（这表示知道性别不会提供关于季节概率的信息，反之亦然；关于独立性，2.5 节会有更多讨论。）

\noindent\textbf{解}\quad 由条件概率定义，
\[
  P(\text{两个都是女孩}|\text{至少一个冬季出生的女孩})
  =
  \frac{P(\text{两个都是女孩，至少一个冬季出生的女孩})}
       {P(\text{至少一个冬季出生的女孩})}.
\]
由于某个特定孩子是冬季出生的女孩的概率是 $1/8$，分母为
\[
  P(\text{至少一个冬季出生的女孩})=1-\left(\frac78\right)^2.
\]
为了计算分子，注意“两个都是女孩，至少一个冬季出生的女孩”与“两个都是女孩，至少一个冬季出生的孩子”是同一事件；再利用性别与季节独立：
\[
\begin{aligned}
  &P(\text{两个都是女孩，至少一个冬季出生的女孩}) \\
  &\quad =
  P(\text{两个都是女孩，至少一个冬季出生的孩子}) \\
  &\quad =
  \left(\frac14\right)P(\text{至少一个冬季出生的孩子}) \\
  &\quad =
  \left(\frac14\right)\left(1-P(\text{两个都不是冬季出生})\right) \\
  &\quad =
  \left(\frac14\right)\left(1-\left(\frac34\right)^2\right).
\end{aligned}
\]
于是
\[
  P(\text{两个都是女孩}|\text{至少一个冬季出生的女孩})
  =
  \frac{(1/4)(1-(3/4)^2)}{1-(7/8)^2}
  =
  \frac{7/64}{15/64}
  =
  \frac{7}{15}.
\]

乍看之下，这个结果显得很荒唐。在例 2.2.5 中，已知至少一个孩子是女孩时，两个孩子都是女孩的条件概率是 $1/3$；为什么知道“至少一个是冬季出生的女孩”后会不同？关键在于，关于出生季节的信息把“至少一个是女孩”推近到了“某个具体孩子是女孩”。以越来越具体的信息为条件，会使概率越来越接近 $1/2$。

例如，以“至少一个孩子是在 3 月 31 日晚上 8:20 出生的女孩”为条件，就非常接近于指定某个孩子；而知道某个具体孩子的信息，并不会给出关于另一个孩子的信息。像出生季节这样看似无关的信息，在例 2.2.5 的两部分之间进行了插值。习题 29 会把这个例子推广到任意与性别独立的特征。

\section*{2.3\quad 贝叶斯公式与全概率公式}

条件概率的定义很简单，只是两个概率的比值，但它有深远后果。第一个后果很容易得到：把定义式中的分母移到等式另一侧即可。

\noindent\textbf{定理 2.3.1（两个事件交集的概率）}\quad 对任意正概率事件 $A$ 与 $B$，
\[
  P(A\cap B)=P(B)P(A|B)=P(A)P(B|A).
\]

这来自对 $P(A|B)$ 的定义两边同乘 $P(B)$，以及对 $P(B|A)$ 的定义两边同乘 $P(A)$。初看起来，这个定理似乎并不十分有用：它只是把条件概率定义稍稍改写了一下，而且用 $P(A|B)$ 来帮助求 $P(A\cap B)$ 似乎有循环之嫌，因为 $P(A|B)$ 正是用 $P(A\cap B)$ 定义的。但我们将看到，这个定理其实非常有用，因为在许多情形下，我们可以不回到定义就直接求条件概率；此时定理 2.3.1 能帮助我们更容易地求 $P(A\cap B)$。

反复应用定理 2.3.1，可以把它推广到 $n$ 个事件的交集。

\noindent\textbf{定理 2.3.2（$n$ 个事件交集的概率）}\quad 对任意事件 $A_1,\ldots,A_n$，若 $P(A_1,A_2,\ldots,A_{n-1})>0$，则
\[
  P(A_1,A_2,\ldots,A_n)
  =
  P(A_1)P(A_2|A_1)P(A_3|A_1,A_2)\cdots
  P(A_n|A_1,\ldots,A_{n-1}).
\]
逗号表示交集，例如 $P(A_3|A_1,A_2)$ 表示 $P(A_3|A_1\cap A_2)$。

事实上，这可以看成 $n!$ 个定理合在一起，因为我们可以任意排列 $A_1,\ldots,A_n$ 而不影响左边。对某些排序，右边往往会比其他排序更容易计算。例如
\[
  P(A_1,A_2,A_3)
  =
  P(A_1)P(A_2|A_1)P(A_3|A_1,A_2)
  =
  P(A_2)P(A_3|A_2)P(A_1|A_2,A_3),
\]
并且还存在另外四种这种形式的展开。要知道该使用哪一种排序，通常需要练习和思考。

现在可以介绍本章的两个主要定理：贝叶斯公式与全概率公式。它们使我们能够在广泛问题中计算条件概率。贝叶斯公式是极其著名且极其有用的结果，它把 $P(A|B)$ 与 $P(B|A)$ 联系起来。

\noindent\textbf{定理 2.3.3（贝叶斯公式）}\quad
\[
  P(A|B)=\frac{P(B|A)P(A)}{P(B)}.
\]

这立即来自定理 2.3.1，而定理 2.3.1 又立即来自条件概率定义。尽管如此，贝叶斯公式在概率与统计中有重要意义和应用，因为我们经常需要求条件概率，而 $P(B|A)$ 往往比 $P(A|B)$ 更容易直接求出（反之亦然）。

贝叶斯公式还可以用\textbf{赔率}而不是概率来写。

\noindent\textbf{定义 2.3.4（赔率）}\quad 事件 $A$ 的赔率为
\[
  \operatorname{odds}(A)=\frac{P(A)}{P(A^c)}.
\]

例如，若 $P(A)=2/3$，我们说支持 $A$ 的赔率是 $2$ 比 $1$（有时写作 $2:1$；也有人说反对 $A$ 的赔率是 $1$ 比 $2$，因此需要小心分辨所说的是支持还是反对某事件的赔率）。当然，也可以从赔率换回概率：
\[
  P(A)=\frac{\operatorname{odds}(A)}{1+\operatorname{odds}(A)}.
\]

把贝叶斯公式中 $P(A|B)$ 的表达式除以 $P(A^c|B)$ 的表达式，就得到贝叶斯公式的赔率形式。

\noindent\textbf{定理 2.3.5（贝叶斯公式的赔率形式）}\quad 对任意正概率事件 $A$ 与 $B$，以 $B$ 为条件后 $A$ 的赔率为
\[
  \frac{P(A|B)}{P(A^c|B)}
  =
  \frac{P(B|A)}{P(B|A^c)}
  \frac{P(A)}{P(A^c)}.
\]

用文字说，就是后验赔率 $P(A|B)/P(A^c|B)$ 等于先验赔率 $P(A)/P(A^c)$ 乘以因子 $P(B|A)/P(B|A^c)$；在统计学中，这个因子称为\textbf{似然比}。有时用贝叶斯公式的这一形式先得到后验赔率会更方便；若需要，再把赔率转回概率。

\textbf{全概率公式}（law of total probability, LOTP）把条件概率与无条件概率联系起来。它对于实现“用条件概率把复杂概率问题分解成较简单片段”这一承诺至关重要，并且经常与贝叶斯公式配合使用。

\noindent\textbf{定理 2.3.6（全概率公式）}\quad 设 $A_1,\ldots,A_n$ 构成样本空间 $S$ 的一个划分（即 $A_i$ 两两不交，且它们的并为 $S$），并且对所有 $i$ 有 $P(A_i)>0$。则
\[
  P(B)=\sum_{i=1}^n P(B|A_i)P(A_i).
\]

\noindent\textbf{证明}\quad 由于 $A_i$ 构成 $S$ 的一个划分，可以把 $B$ 分解为
\[
  B=(B\cap A_1)\cup(B\cap A_2)\cup\cdots\cup(B\cap A_n).
\]
这在图 2.3 中有所说明：我们把 $B$ 切成从 $B\cap A_1$ 到 $B\cap A_n$ 的若干小块。由概率第二公理，这些小块两两不交，因此可以把它们的概率相加得到 $P(B)$：
\[
  P(B)=P(B\cap A_1)+P(B\cap A_2)+\cdots+P(B\cap A_n).
\]
现在对每个 $P(B\cap A_i)$ 应用定理 2.3.1：
\[
  P(B)=P(B|A_1)P(A_1)+\cdots+P(B|A_n)P(A_n).
\]

全概率公式告诉我们：为了得到 $B$ 的无条件概率，可以把样本空间划分成互不相交的片段 $A_i$，求出 $B$ 在每个片段内的条件概率，然后对这些条件概率作加权和，权重为 $P(A_i)$。如何划分样本空间至关重要：选得好的划分会把复杂问题化为较简单的片段；选得差的划分只会加剧问题，迫使我们计算 $n$ 个困难概率，而不是只计算一个。

接下来的几个例子展示如何把贝叶斯公式与全概率公式结合起来，根据观察到的证据更新我们的信念。

\begin{figure}[htbp]
  \centering
  \includegraphics[width=0.66\textwidth]{2-3.png}
  \caption{$A_i$ 构成样本空间的划分；$P(B)$ 等于 $\sum_i P(B\cap A_i)$。}
\end{figure}

\noindent\textbf{例 2.3.7（随机硬币）}\quad 你有一枚公平硬币，和一枚偏硬币；偏硬币落到正面的概率为 $3/4$。随机选取其中一枚硬币并抛掷三次，结果三次都是正面。在这一信息下，你选到公平硬币的概率是多少？

\noindent\textbf{解}\quad 令 $A$ 为选中的硬币三次都落到正面的事件，令 $F$ 为选到公平硬币的事件。我们关心 $P(F|A)$，但由于知道是哪枚硬币有助于判断三次正面的概率，所以 $P(A|F)$ 与 $P(A|F^c)$ 更容易求；这提示我们使用贝叶斯公式与全概率公式。于是
\[
\begin{aligned}
  P(F|A)
  &=
  \frac{P(A|F)P(F)}{P(A)} \\
  &=
  \frac{P(A|F)P(F)}
       {P(A|F)P(F)+P(A|F^c)P(F^c)} \\
  &=
  \frac{(1/2)^3\cdot 1/2}
       {(1/2)^3\cdot 1/2+(3/4)^3\cdot 1/2} \\
  &\approx 0.23.
\end{aligned}
\]
在抛掷硬币之前，我们认为选到公平硬币和偏硬币等可能：$P(F)=P(F^c)=1/2$。然而，观察到三次正面之后，选到偏硬币就比选到公平硬币更可能，所以 $P(F|A)$ 只有约 $0.23$。

\noindent\textbf{警示 2.3.8（先验与后验）}\quad 在上例计算中，第一步之后若说“$P(A)=1$，因为我们知道 $A$ 已经发生”，这是\textit{不正确}的。的确，$P(A|A)=1$，但 $P(A)$ 是 $A$ 的先验概率，$P(F)$ 是 $F$ 的先验概率；二者都是在观察实验中的任何数据之前的概率。它们不能与以证据 $A$ 为条件的后验概率混淆。

\noindent\textbf{例 2.3.9（罕见病检测）}\quad 一位名叫 Fred 的患者接受一种名为 conditionitis 的疾病检测；这种病影响 $1\%$ 的人群。检测结果为阳性，也就是说，检测声称 Fred 患有该病。令 $D$ 表示 Fred 患病的事件，令 $T$ 表示检测为阳性的事件。

假设该检测“$95\%$ 准确”。检测准确性有不同度量，但在本题中假设它表示 $P(T|D)=0.95$ 且 $P(T^c|D^c)=0.95$。量 $P(T|D)$ 称为检测的\textbf{灵敏度}或\textbf{真阳性率}，而 $P(T^c|D^c)$ 称为\textbf{特异度}或\textbf{真阴性率}。

求在检测结果提供的证据下，Fred 患有 conditionitis 的条件概率。

\noindent\textbf{解}\quad 应用贝叶斯公式与全概率公式，
\[
\begin{aligned}
  P(D|T)
  &=
  \frac{P(T|D)P(D)}{P(T)} \\
  &=
  \frac{P(T|D)P(D)}
       {P(T|D)P(D)+P(T|D^c)P(D^c)} \\
  &=
  \frac{0.95\cdot 0.01}
       {0.95\cdot 0.01+0.05\cdot 0.99} \\
  &\approx 0.16.
\end{aligned}
\]
因此，即使检测看起来相当可靠，在 Fred 检测呈阳性的条件下，他真正患有 conditionitis 的概率也只有 $16\%$。

许多人会惊讶地发现，尽管检测有 $95\%$ 的准确率，给定阳性结果后患病的条件概率却只有 $16\%$。理解这个很低的后验患病概率的关键，是认识到这里有两个因素同时起作用：检测证据，以及我们关于疾病流行率的\textit{先验信息}。

尽管检测提供了支持患病的证据，conditionitis 仍然是一种罕见病。条件概率 $P(D|T)$ 反映的是这两个因素之间的平衡：它恰当地把疾病的罕见性与错误检测结果的罕见性相权衡。

为了获得进一步直觉，考虑图 2.4 所示的 $10000$ 人总体，其中 $100$ 人患有 conditionitis，$9900$ 人没有；这对应于 $1\%$ 的患病率。如果对总体中所有人进行检测，那么在 $100$ 名患病者中，预期 $95$ 人检测阳性、$5$ 人检测阴性。在 $9900$ 名健康者中，预期 $(0.95)(9900)\approx 9405$ 人检测阴性，$495$ 人检测阳性。

\begin{figure}[htbp]
  \centering
  \includegraphics[width=0.62\textwidth]{2-4.png}
  \caption{在 $10000$ 人总体中检测罕见病，患病率为 $1\%$，真阳性率与真阴性率均为 $95\%$。气泡大小不按比例绘制。}
\end{figure}

现在聚焦于检测为阳性的个体，也就是以阳性检测结果为条件。$95$ 个真阳性（检测阳性且患病者）远少于 $495$ 个假阳性（未患病但检测阳性者）。因此，检测阳性的人大多数其实并没有患病。

\noindent\textbf{例 2.3.10（六指男子）}\quad 某国发生了一起犯罪。罪犯是该国 $n$ 名男子中的一个且只有一个。初始时，这 $n$ 名男子都被认为等可能是罪犯。随后，一名目击者报告说，犯罪者是一名右手有六根手指的男子。

令 $p_0$ 为无辜男子右手有六根手指的概率，令 $p_1$ 为罪犯右手有六根手指的概率，并且 $p_0<p_1$。（可能有 $p_1<1$，因为目击者并非 $100\%$ 可靠。）令
\[
  a=\frac{p_0}{p_1},\qquad b=\frac{1-p_1}{1-p_0}.
\]
Rugen 住在该国，并被发现右手有六根手指。

\begin{enumerate}
  \item[(a)] 给定这一信息，Rugen 是罪犯的概率是多少？
  \item[(b)] 现在假设该国所有 $n$ 名男子都检查了双手，并且 Rugen 是唯一右手有六根手指的人。给定这一信息，Rugen 是罪犯的概率是多少？
\end{enumerate}

\noindent\textbf{解}\quad (a) 令 $R$ 为 Rugen 有罪的事件，令 $M$ 为他右手有六根手指的事件。由贝叶斯公式与全概率公式，
\[
  P(R|M)
  =
  \frac{P(M|R)P(R)}
       {P(M|R)P(R)+P(M|R^c)P(R^c)}
  =
  \frac{p_1\cdot \frac1n}
       {p_1\cdot \frac1n+p_0(1-\frac1n)}
  =
  \frac{1}{1+a(n-1)}.
\]

(b) 令 $N$ 为该国除 Rugen 外没有男子右手有六根手指的事件。沿用上面的记号，
\[
\begin{aligned}
  P(R|M,N)
  &=
  \frac{P(M,N|R)P(R)}
       {P(M,N|R)P(R)+P(M,N|R^c)P(R^c)} \\
  &=
  \frac{p_1(1-p_0)^{n-1}\cdot \frac1n}
       {p_1(1-p_0)^{n-1}\cdot \frac1n
        +p_0(1-p_1)(1-p_0)^{n-2}(1-\frac1n)} \\
  &=
  \frac{1}{1+ab(n-1)}.
\end{aligned}
\]

\section*{2.4\quad 条件概率仍是概率}

当我们以事件 $E$ 为条件时，我们是在更新信念，使其与这一知识一致；这等于把自己放进一个知道 $E$ 已经发生的宇宙中。不过，在这个新的宇宙里，概率法则照常运作。条件概率满足概率的所有性质。因此，我们关于概率推导出的任何结果，在把所有无条件概率替换成以 $E$ 为条件的概率之后，仍然有效。特别地：

\begin{itemize}
  \item 条件概率介于 $0$ 与 $1$ 之间。
  \item $P(S|E)=1,\quad P(\varnothing|E)=0$。
  \item 若 $A_1,A_2,\ldots$ 两两不交，则
  \[
    P\left(\bigcup_{j=1}^{\infty}A_j\,\middle|\,E\right)
    =
    \sum_{j=1}^{\infty}P(A_j|E).
  \]
  \item $P(A^c|E)=1-P(A|E)$。
  \item 容斥公式：
  \[
    P(A\cup B|E)=P(A|E)+P(B|E)-P(A\cap B|E).
  \]
\end{itemize}

\noindent\textbf{警示 2.4.1}\quad 写 $P(A|E)$ 时，并不表示 $A|E$ 是一个事件、而我们正在取它的概率；$A|E$ 不是事件。相反，$P(\cdot|E)$ 是一个概率函数，它根据 $E$ 已经发生这一知识分配概率；而 $P(\cdot)$ 是另一个概率函数，它分配概率时不考虑 $E$ 是否已经发生。把事件 $A$ 代入 $P(\cdot)$ 会得到一个数 $P(A)$；把它代入 $P(\cdot|E)$ 则会得到另一个数 $P(A|E)$，其中纳入了知道 $E$ 已发生所提供的信息（如果有的话）。

为了从数学上证明条件概率确实是概率，固定一个满足 $P(E)>0$ 的事件 $E$，并对任意事件 $A$ 定义
\[
  \widetilde P(A)=P(A|E).
\]
这个记号强调：我们固定 $E$，并把 $P(\cdot|E)$ 当作新的概率函数。我们只需检查概率的两条公理。首先，
\[
  \widetilde P(\varnothing)=P(\varnothing|E)
  =
  \frac{P(\varnothing\cap E)}{P(E)}=0,\qquad
  \widetilde P(S)=P(S|E)
  =
  \frac{P(S\cap E)}{P(E)}=1.
\]
其次，若 $A_1,A_2,\ldots$ 两两不交，则
\[
  \widetilde P(A_1\cup A_2\cup\cdots)
  =
  \frac{P((A_1\cap E)\cup(A_2\cap E)\cup\cdots)}{P(E)}
  =
  \frac{\sum_{j=1}^{\infty}P(A_j\cap E)}{P(E)}
  =
  \sum_{j=1}^{\infty}\widetilde P(A_j).
\]
因此，$\widetilde P$ 满足概率公理。

反过来，\textit{所有}概率都可以看成条件概率：每当我们作出一个概率陈述时，总有一些背景信息是我们正在条件化的，即使没有明确说明。考虑本章开头的下雨例子。自然的做法是把初始概率建立在
过去下雨天数的比例上。但我们该看过去的哪些天？如果今天是 11 月 1 日，是否只应计算过去秋天里的雨天，从而以季节为条件？是否还应以精确月份或精确日期为条件？关于地点也可以提出同样的问题：我们是否只看精确地点曾经下雨的日子，还是附近某处下过雨也足够？

为了确定看似无条件的概率 $P(R)$，我们实际上必须决定要以哪些背景信息为条件。这些选择需要仔细思考，不同的人可能给出不同的先验概率 $P(R)$（尽管每个人都可以同意如何根据新证据进行更新）。

既然所有概率都以背景信息为条件，我们可以想象总有一条条件竖线，而背景知识 $K$ 位于竖线右侧。那么无条件概率 $P(A)$ 只是 $P(A|K)$ 的简写；背景知识被吸收到字母 $P$ 中，而没有显式写出来。

把上面的讨论压缩成一句话：
\[
  \text{\textit{条件概率是概率，而所有概率都是条件的。}}
\]

现在给出贝叶斯公式和全概率公式的条件形式。它们通过在普通的贝叶斯公式与全概率公式中，把 $E$ 加到每一处条件竖线右侧而得到。

\noindent\textbf{定理 2.4.2（带额外条件的贝叶斯公式）}\quad 若 $P(A\cap E)>0$ 且 $P(B\cap E)>0$，则
\[
  P(A|B,E)=\frac{P(B|A,E)P(A|E)}{P(B|E)}.
\]

\noindent\textbf{定理 2.4.3（带额外条件的全概率公式）}\quad 设 $A_1,\ldots,A_n$ 构成 $S$ 的一个划分。若对所有 $i$ 有 $P(A_i\cap E)>0$，则
\[
  P(B|E)=\sum_{i=1}^n P(B|A_i,E)P(A_i|E).
\]

贝叶斯公式和全概率公式的这些额外条件形式，可以像我们验证 $\widetilde P$ 满足概率公理那样证明；它们也直接来自“条件概率是概率”这一元定理。

\noindent\textbf{例 2.4.4（随机硬币，续）}\quad 延续例 2.3.7 的情形，假设现在已经看到所选硬币连续三次落到正面。如果第四次抛掷这枚硬币，它再次落到正面的概率是多少？

\noindent\textbf{解}\quad 和以前一样，令 $A$ 为所选硬币前三次都落到正面的事件，并定义新事件 $H$ 表示所选硬币第四次落到正面。我们关心 $P(H|A)$。若知道手中是否为公平硬币，会非常有帮助。带额外条件的全概率公式给出 $P(H|A)$，它是 $P(H|F,A)$ 与 $P(H|F^c,A)$ 的加权平均；而在这两个条件概率中，我们确实知道自己是否拿到公平硬币：
\[
\begin{aligned}
  P(H|A)
  &=
  P(H|F,A)P(F|A)+P(H|F^c,A)P(F^c|A)\\
  &\approx
  \frac12\cdot 0.23+\frac34\cdot(1-0.23)\\
  &\approx 0.69.
\end{aligned}
\]
这里的后验概率 $P(F|A)$ 与 $P(F^c|A)$ 来自例 2.3.7 的答案。

等价的解法是定义一个新的概率函数 $\widetilde P$，使得对任意事件 $B$ 有 $\widetilde P(B)=P(B|A)$。这个新函数分配的是已经根据 $A$ 发生这一知识更新后的概率。于是由普通全概率公式，
\[
  \widetilde P(H)=\widetilde P(H|F)\widetilde P(F)+\widetilde P(H|F^c)\widetilde P(F^c),
\]
这与使用带额外条件的全概率公式完全相同。这再次说明：条件概率是概率。

\noindent\textbf{例 2.4.5（全体一致）}\quad Lisa Zyga 的文章《为什么太多证据可能是坏事》提到：
\begin{quote}
在古代犹太法律中，如果受审嫌疑人被所有法官一致判定有罪，那么嫌疑人会被宣告无罪。这个推理听起来违反直觉，但当时的立法者注意到，全体一致往往表示司法过程中存在系统性错误。
\end{quote}

有 $n$ 名法官裁决一个案件。嫌疑人有罪的先验概率为 $p$。每位法官投票决定定罪或无罪。以概率 $s$ 会发生系统性错误（例如辩方无能）。若系统性错误发生，则法官会一致投票定罪，也就是所有 $n$ 名法官都投定罪票。

系统性错误是否发生，与嫌疑人是否有罪独立。在没有系统性错误且嫌疑人有罪的条件下，每位法官独立地以概率 $c$ 投票定罪。在没有系统性错误且嫌疑人无罪的条件下，每位法官独立地以概率 $w$ 投票定罪。假设
\[
  0<p<1,\qquad 0<s<1,\qquad 0<w<\frac12<c<1.
\]

\begin{enumerate}
  \item[(a)] 仅在本问中，假设 $n$ 名法官中恰有 $k$ 名投票定罪，其中 $k<n$。给定这一信息，求嫌疑人有罪的概率。
  \item[(b)] 现在假设所有 $n$ 名法官都投票定罪。给定这一信息，求嫌疑人有罪的概率。
  \item[(c)] 把 (b) 的答案看成 $n$ 的函数，它是递增函数吗？用文字给出一个简短的直观解释。
\end{enumerate}

\noindent\textbf{解}\quad (a) 由于 $k<n$，系统性错误没有发生。本问中我们隐含以这一点为条件。令 $G$ 为嫌疑人有罪的事件，令 $X$ 为投票定罪的法官人数。使用贝叶斯公式、全概率公式和二项分布概率质量函数，
\[
  P(G|X=k)
  =
  \frac{P(X=k|G)P(G)}{P(X=k)}
  =
  \frac{pc^k(1-c)^{n-k}}
       {pc^k(1-c)^{n-k}+(1-p)w^k(1-w)^{n-k}}.
\]

(b) 令 $U$ 为事件 $X=n$，令 $B$ 为发生系统性错误的事件。则
\[
  P(G|U)
  =
  \frac{P(U|G)P(G)}{P(U)}
  =
  \frac{pP(U|G)}
       {pP(U|G)+(1-p)P(U|G^c)}.
\]
由带额外条件的全概率公式，
\[
  P(U|G)=P(U|G,B)P(B|G)+P(U|G,B^c)P(B^c|G)
  =
  s+(1-s)c^n,
\]
\[
  P(U|G^c)=P(U|G^c,B)P(B|G^c)+P(U|G^c,B^c)P(B^c|G^c)
  =
  s+(1-s)w^n.
\]
因此
\[
  P(G|U)
  =
  \frac{p(s+(1-s)c^n)}
       {p(s+(1-s)c^n)+(1-p)(s+(1-s)w^n)}.
\]

(c) 不是。因为较大的 $n$ 会带来较高的系统性错误可能性；一旦发生系统性错误，法官们的投票对嫌疑人是否有罪就不提供信息。当 $n\to\infty$ 时，(b) 的答案会回到先验概率 $p$。

我们常常想以不止一条信息为条件，而现在已有若干方法可以做到。例如，为了求 $P(A|B,C)$，可以采用下面几种思路。

\begin{enumerate}
  \item 把 $B,C$ 看成单个事件 $B\cap C$，并使用条件概率定义：
  \[
    P(A|B,C)=\frac{P(A,B,C)}{P(B,C)}.
  \]
  如果最容易把 $B$ 与 $C$ 联合起来思考，这就是很自然的方法。随后可以尝试计算分子和分母。例如，可以在分子和分母中都使用全概率公式；或者把分子写成 $P(B,C|A)P(A)$（这会给出贝叶斯公式的一种形式），并用全概率公式帮助计算分母。

  \item 使用以 $C$ 为额外条件的贝叶斯公式：
  \[
    P(A|B,C)=\frac{P(B|A,C)P(A|C)}{P(B|C)}.
  \]
  如果我们想把问题中的一切都看成以 $C$ 为条件，这就是自然的方法。

  \item 使用以 $B$ 为额外条件的贝叶斯公式：
  \[
    P(A|B,C)=\frac{P(C|A,B)P(A|B)}{P(C|B)}.
  \]
  这与前一种方法相同，只是交换了 $B$ 与 $C$ 的角色。单独提到它，是为了强调：不思考哪个事件应扮演哪个角色就把量塞进公式，是不好的做法。
\end{enumerate}

处理这类条件问题有多种方法，这一点既有挑战性，也很有力量。

\section*{2.5\quad 事件的独立性}

我们已经看到许多例子：以一个事件为条件会改变我们对另一个事件概率的信念。若事件彼此不提供信息，则称为\textbf{独立}。

\noindent\textbf{定义 2.5.1（两个事件的独立性）}\quad 若事件 $A$ 与 $B$ 满足
\[
  P(A\cap B)=P(A)P(B),
\]
则称 $A$ 与 $B$ 独立。若 $P(A)>0$ 且 $P(B)>0$，这等价于
\[
  P(A|B)=P(A),
\]
也等价于 $P(B|A)=P(B)$。

用文字说，两个事件独立，是指它们交集的概率可以由各自概率相乘得到。等价地，若知道 $B$ 已经发生不会改变我们对 $A$ 发生概率的判断（反之亦然），则 $A$ 与 $B$ 独立。注意，独立性是一种对称关系：如果 $A$ 独立于 $B$，那么 $B$ 也独立于 $A$。

\noindent\textbf{警示 2.5.2}\quad 独立性与互不相交完全不同。若 $A$ 与 $B$ 互不相交，则 $P(A\cap B)=0$，所以互不相交的事件只有在 $P(A)=0$ 或 $P(B)=0$ 时才可能独立。知道 $A$ 发生告诉我们 $B$ 一定没有发生，因此 $A$ 显然传递了关于 $B$ 的信息，意味着二者并不独立（除非 $A$ 或 $B$ 已经具有零概率）。

直观上，如果 $A$ 不提供关于 $B$ 是否发生的信息，那么它也不应提供关于 $B^c$ 是否发生的信息。下面证明一个有用的结果。

\noindent\textbf{命题 2.5.3}\quad 若 $A$ 与 $B$ 独立，则 $A$ 与 $B^c$ 独立，$A^c$ 与 $B$ 独立，且 $A^c$ 与 $B^c$ 独立。

\noindent\textbf{证明}\quad 设 $A$ 与 $B$ 独立。先证明 $A$ 与 $B^c$ 独立。若 $P(A)=0$，则 $A$ 与包括 $B^c$ 在内的所有事件独立。故假设 $P(A)\ne 0$。于是
\[
  P(B^c|A)=1-P(B|A)=1-P(B)=P(B^c),
\]
所以 $A$ 与 $B^c$ 独立。交换 $A$ 与 $B$ 的角色，可知 $A^c$ 与 $B$ 独立。再利用 $A,B$ 独立推出 $A,B^c$ 独立，并让 $A^c$ 扮演 $A$ 的角色，也得到 $A^c$ 与 $B^c$ 独立。

我们也常需要讨论三个或更多事件的独立性。

\noindent\textbf{定义 2.5.4（三个事件的独立性）}\quad 若事件 $A,B,C$ 满足以下所有等式：
\[
  P(A\cap B)=P(A)P(B),
\]
\[
  P(A\cap C)=P(A)P(C),
\]
\[
  P(B\cap C)=P(B)P(C),
\]
\[
  P(A\cap B\cap C)=P(A)P(B)P(C),
\]
则称 $A,B,C$ 独立。若前三个条件成立，则称 $A,B,C$ \textbf{两两独立}。两两独立并不推出独立：可能只知道 $A$ 或只知道 $B$ 对预测 $C$ 是否发生没有帮助，但同时知道 $A$ 与 $B$ 都发生，却仍然与 $C$ 高度相关。下面是一个简单例子。

\noindent\textbf{例 2.5.5（两两独立不推出独立）}\quad 考虑两次公平且相互独立的抛硬币。令 $A$ 为第一次为正面的事件，$B$ 为第二次为正面的事件，$C$ 为两次结果相同的事件。则 $A,B,C$ 两两独立但不独立，因为
\[
  P(A\cap B\cap C)=\frac14,\qquad P(A)P(B)P(C)=\frac18.
\]
关键在于，只知道 $A$ 或只知道 $B$ 并不能告诉我们关于 $C$ 的信息；但知道 $A$ 与 $B$ 都发生，就给出了关于 $C$ 的信息（事实上，在这个例子中给出了关于 $C$ 的完全信息）。

另一方面，$P(A\cap B\cap C)=P(A)P(B)P(C)$ 并不推出两两独立；这可以通过极端情形 $P(A)=0$ 快速看出，此时等式变为 $0=0$，它没有告诉我们关于 $B$ 与 $C$ 的任何信息。

可以类似地定义任意多个事件的独立性。直观上，思想是：知道这些事件中任意特定子集发生了什么，都不会提供关于不在该子集中的事件发生情况的信息。

\noindent\textbf{定义 2.5.6（多个事件的独立性）}\quad 对 $n$ 个事件 $A_1,A_2,\ldots,A_n$，若要称它们独立，需要任意一对满足
\[
  P(A_i\cap A_j)=P(A_i)P(A_j)\quad (i\ne j),
\]
任意三元组满足
\[
  P(A_i\cap A_j\cap A_k)=P(A_i)P(A_j)P(A_k)
\]
（其中 $i,j,k$ 互不相同），并且任意四元组、五元组等也都满足类似等式。这很快会变得笨重；后面我们会讨论思考独立性的其他方法。对于无限多个事件，若每个有限子集都是独立的，则称这些事件独立。

条件独立性的定义与独立性类似。

\noindent\textbf{定义 2.5.7（条件独立性）}\quad 若事件 $A$ 与 $B$ 满足
\[
  P(A\cap B|E)=P(A|E)P(B|E),
\]
则称 $A$ 与 $B$ 在给定 $E$ 时\textbf{条件独立}。

\noindent\textbf{警示 2.5.8}\quad 混淆独立性与条件独立性很容易造成严重错误。两个事件可以在给定 $E$ 时条件独立，却在给定 $E^c$ 时不条件独立。两个事件可以在给定 $E$ 时条件独立，但并不独立。两个事件也可以独立，但在给定 $E$ 时不条件独立。

特别地，$P(A,B)=P(A)P(B)$ 并不推出
\[
  P(A,B|E)=P(A|E)P(B|E).
\]
我们不能像从全概率公式得到带额外条件的全概率公式那样，只是到处插入“给定 $E$”。原因在于，全概率公式总是成立（它是概率公理的后果），而 $P(A,B)$ 是否等于 $P(A)P(B)$ 则取决于 $A$ 与 $B$。

接下来的几个例子说明这些区别。处理条件概率和条件独立性时必须非常谨慎。

\noindent\textbf{例 2.5.9（给定 $E$ 与给定 $E^c$ 的条件独立）}\quad 假设课程分为两类：好课和坏课。在好课中，如果你努力学习，就很可能得到 A；在坏课中，教授不管学生努力程度如何都会随机给分。令 $G$ 为课程是好课的事件，$W$ 为你努力学习的事件，$A$ 为你得到 A 的事件。则 $W$ 与 $A$ 在给定 $G^c$ 时条件独立，但它们在给定 $G$ 时不条件独立。

\noindent\textbf{例 2.5.10（条件独立不推出独立）}\quad 再回到例 2.3.7 的情形。假设我们选中的是公平硬币，或是一枚正面概率为 $3/4$ 的偏硬币，但不知道选中了哪一枚。我们抛掷这枚硬币若干次。给定选中了公平硬币，抛掷结果彼此独立，每次正面概率为 $1/2$。类似地，给定选中了偏硬币，抛掷结果也彼此独立，每次正面概率为 $3/4$。

然而，抛掷结果并不是无条件独立的，因为如果不知道选中了哪枚硬币，那么观察到一串抛掷结果会提供关于手中是公平硬币还是偏硬币的信息；这反过来有助于预测同一枚硬币未来的抛掷结果。

形式化地说，令 $F$ 为选中了公平硬币的事件，令 $A_1,A_2$ 分别为第一次、第二次抛掷为正面的事件。给定 $F$ 时，$A_1$ 与 $A_2$ 独立；但 $A_1$ 与 $A_2$ 并非无条件独立，因为 $A_1$ 提供了关于 $A_2$ 的信息。

\noindent\textbf{例 2.5.11（独立不推出条件独立）}\quad 我的朋友 Alice 和 Bob 是唯二会给我打电话的人。每天，他们独立地决定当天是否给我打电话。令 $A$ 为 Alice 下周五给我打电话的事件，$B$ 为 Bob 下周五给我打电话的事件。假设 $A$ 与 $B$ 无条件独立，且 $P(A)>0,P(B)>0$。

然而，给定我下周五恰好接到一个电话，$A$ 与 $B$ 就不再独立：这个电话来自 Alice 当且仅当它不来自 Bob。换言之，令 $C$ 为我下周五恰好接到一个电话的事件，则 $P(B|C)>0$，但 $P(B|A,C)=0$，所以 $A$ 与 $B$ 在给定 $C$ 时不条件独立。

\noindent\textbf{例 2.5.12（婴儿为什么哭？）}\quad 某婴儿哭泣当且仅当她饿了、累了，或二者兼有。令 $C$ 为婴儿哭的事件，$H$ 为她饿了的事件，$T$ 为她累了的事件。令 $P(C)=c,\ P(H)=h,\ P(T)=t$，其中 $c,h,t$ 均不等于 $0$ 或 $1$。设 $H$ 与 $T$ 独立。

\begin{enumerate}
  \item[(a)] 用 $h,t$ 表示 $c$。
  \item[(b)] 求 $P(H|C),P(T|C)$ 和 $P(H,T|C)$。
  \item[(c)] $H$ 与 $T$ 在给定 $C$ 时条件独立吗？请用两种方式解释：一种是利用 (b) 中的量进行代数说明，另一种用文字给出直观说明。
\end{enumerate}

\noindent\textbf{解}\quad (a) 由于 $H$ 与 $T$ 独立，
\[
  P(C)=P(H\cup T)=P(H)+P(T)-P(H\cap T)=h+t-ht.
\]

(b) 由贝叶斯公式，
\[
  P(H|C)=\frac{P(C|H)P(H)}{P(C)}=\frac{h}{c},
\]
\[
  P(T|C)=\frac{P(C|T)P(T)}{P(C)}=\frac{t}{c},
\]
\[
  P(H,T|C)=\frac{P(C|H,T)P(H,T)}{P(C)}=\frac{ht}{c}.
\]

(c) 不，$H$ 与 $T$ 在给定 $C$ 时不条件独立，因为
\[
  P(H,T|C)=\frac{ht}{c}<\frac{ht}{c^2}=P(H|C)P(T|C).
\]
直观上也可以看出它们为何在给定 $C$ 时不条件独立：若婴儿正在哭，但并不饿，那她一定是累了。

\section*{2.6\quad 贝叶斯公式的一致性}

贝叶斯公式的一个重要性质是它具有\textbf{一致性}：如果我们接收到多条信息，并希望更新概率以纳入所有信息，那么逐条更新、一次只考虑一条证据，和同时使用所有证据一次性更新，结果没有区别。比如，假设我们正在进行一个为期一周的实验，每天结束时都会产生数据。我们可以每天使用贝叶斯公式根据当天数据更新概率；也可以整周去度假，周五下午回来后用整周数据一次性更新。两种方法会给出相同结果。

下面看这一原则的一个具体应用。

\noindent\textbf{例 2.6.1（罕见病检测，续）}\quad Fred 在例 2.3.9 中检测 conditionitis 为阳性；他决定再检测一次。新的检测在给定患病状态后与原检测独立，并且具有相同的灵敏度和特异度。不幸的是，Fred 第二次检测也呈阳性。请用两种方法求给定证据下 Fred 患病的概率：一种是一步法，同时以两次检测结果为条件；另一种是两步法，先根据第一次检测更新概率，再根据第二次检测继续更新。

\noindent\textbf{解}\quad 令 $D$ 为 Fred 患病的事件，$T_1$ 为第一次检测为阳性的事件，$T_2$ 为第二次检测为阳性的事件。在例 2.3.9 中，我们用贝叶斯公式和全概率公式求出了 $P(D|T_1)$。另一种快速解法是使用贝叶斯公式的赔率形式：
\[
  \frac{P(D|T_1)}{P(D^c|T_1)}
  =
  \frac{P(D)}{P(D^c)}\frac{P(T_1|D)}{P(T_1|D^c)}
  =
  \frac{1}{99}\cdot\frac{0.95}{0.05}
  \approx 0.19.
\]
由于 $P(D|T_1)/(1-P(D|T_1))=0.19$，可得
\[
  P(D|T_1)=\frac{0.19}{1+0.19}\approx 0.16,
\]
与先前答案一致。在本例中，贝叶斯公式的赔率形式更快，因为我们不必计算普通贝叶斯公式分母中的无条件概率 $P(T_1)$。现在继续用赔率形式，看看第二次检测为阳性后会发生什么。

\noindent\textit{一步法：}同时根据两次检测结果更新，
\[
\begin{aligned}
  \frac{P(D|T_1\cap T_2)}{P(D^c|T_1\cap T_2)}
  &=
  \frac{P(D)}{P(D^c)}
  \frac{P(T_1\cap T_2|D)}{P(T_1\cap T_2|D^c)}\\
  &=
  \frac{1}{99}\cdot\frac{0.95^2}{0.05^2}
  =
  \frac{361}{99}
  \approx 3.646,
\end{aligned}
\]
这对应于概率约为 $0.78$。

\noindent\textit{两步法：}第一次检测之后，Fred 患病的后验赔率为
\[
  \frac{P(D|T_1)}{P(D^c|T_1)}
  =
  \frac{1}{99}\cdot\frac{0.95}{0.05}
  \approx 0.19.
\]
这个后验赔率成为新的先验赔率；随后根据第二次检测更新，得到
\[
\begin{aligned}
  \frac{P(D|T_1\cap T_2)}{P(D^c|T_1\cap T_2)}
  &=
  \frac{P(D|T_1)}{P(D^c|T_1)}
  \frac{P(T_2|D,T_1)}{P(T_2|D^c,T_1)}\\
  &=
  \left(\frac{1}{99}\cdot\frac{0.95}{0.05}\right)
  \frac{0.95}{0.05}
  =
  \frac{361}{99}
  \approx 3.646,
\end{aligned}
\]
与上面相同。

注意，第二次阳性检测结果使 Fred 患病的概率从 $0.16$ 跃升到 $0.78$，使我们更有把握认为 Fred 确实患有 conditionitis。这个故事的寓意是：听取第二意见是个好主意。

\section*{2.7\quad 条件化作为解题工具}

条件化是强大的解题工具，因为它让我们可以进行“愿望式思考”：当遇到某个问题，而如果知道事件 $E$ 是否发生就会使问题变容易时，我们可以先以 $E$ 为条件，再以 $E^c$ 为条件，分别考虑这两种可能，最后用全概率公式合并。

\subsection*{2.7.1\quad 策略：以你希望知道的信息为条件}

\noindent\textbf{例 2.7.1（Monty Hall）}\quad 在 Monty Hall 主持的电视节目《Let's Make a Deal》中，参赛者从三扇关闭的门中选择一扇；其中两扇门后是山羊，一扇门后是汽车。知道汽车位置的 Monty 随后在剩余两扇门中打开一扇。被他打开的门后总是山羊（他绝不会展示汽车）。若他有选择，则等可能地随机选择一扇门打开。随后 Monty 让参赛者选择是否换到另一扇未打开的门。若参赛者的目标是得到汽车，她是否应该换门？

\begin{center}
  \includegraphics[width=0.42\textwidth]{2-4-1.png}
\end{center}

\noindent\textbf{解}\quad 将门标记为 $1,2,3$。不失一般性，假设参赛者选择了门 $1$（若她没有选择门 $1$，我们可以重新标记门，或把下面的解答中的门号相应置换）。Monty 打开一扇门，展示一只山羊。当参赛者决定是否换到另一扇未打开的门时，她真正希望知道什么？自然地，如果她知道汽车在哪里，决定就会容易得多。这提示我们应以汽车的位置为条件。令 $C_i$ 为汽车在门 $i$ 后的事件，$i=1,2,3$。由全概率公式，
\[
  P(\text{得到汽车})
  =
  P(\text{得到汽车}|C_1)\cdot\frac13
  +P(\text{得到汽车}|C_2)\cdot\frac13
  +P(\text{得到汽车}|C_3)\cdot\frac13.
\]

假设参赛者采用换门策略。若汽车在门 $1$ 后，则换门会失败，所以 $P(\text{得到汽车}|C_1)=0$。若汽车在门 $2$ 或门 $3$ 后，因为 Monty 总会展示一只山羊，剩下未打开的门必然含有汽车，所以换门会成功。因此
\[
  P(\text{得到汽车})
  =
  0\cdot\frac13+1\cdot\frac13+1\cdot\frac13
  =
  \frac23.
\]
所以换门策略有 $2/3$ 的成功率，参赛者应该换到另一扇门。

\begin{figure}[htbp]
  \centering
  \includegraphics[width=0.72\textwidth]{2-5.png}
  \caption{Monty Hall 问题的树图。采用换门策略时，参赛者在汽车位于门 $2$ 或门 $3$ 后时获胜，因此成功概率为 $2/3$。}
\end{figure}

图 2.5 是刚才论证的树图：在换门策略下，只要汽车在门 $2$ 或门 $3$ 后，参赛者就会获胜，而这件事的概率为 $2/3$。也可以给出支持换门的频率学派直觉论证。设想玩这个游戏 $1000$ 次。通常大约有 $333$ 次你最初猜中汽车位置，此时换门会失败；另外约 $667$ 次，你会通过换门获胜。

不过这里有一个细微之处：参赛者决定是否换门时，也知道了 Monty 打开的是哪扇门。我们已经说明，在采用换门策略时，成功的无条件概率是 $2/3$；但也应说明，在给定 Monty 所提供信息后，换门成功的条件概率仍为 $2/3$。

令 $M_j$ 为 Monty 打开门 $j$ 的事件，其中 $j=2,3$。则
\[
  P(\text{得到汽车})
  =
  P(\text{得到汽车}|M_2)P(M_2)
  +
  P(\text{得到汽车}|M_3)P(M_3),
\]
其中由对称性，$P(M_2)=P(M_3)=1/2$，且
\[
  P(\text{得到汽车}|M_2)=P(\text{得到汽车}|M_3).
\]
这里的对称性在于，题目叙述中没有任何东西区分门 $2$ 与门 $3$；相反，习题 40 会考虑 Monty 更喜欢打开门 $2$ 而不是门 $3$ 的情形。

令
\[
  x=P(\text{得到汽车}|M_2)=P(\text{得到汽车}|M_3).
\]
代入已知量得
\[
  \frac23=P(\text{得到汽车})=\frac{x}{2}+\frac{x}{2}=x,
\]
如所声称。

贝叶斯公式也可以很好地用于在给定证据后求换门策略成功的条件概率。假设 Monty 打开门 $2$。沿用上面的记号与结果，
\[
  P(C_1|M_2)=
  \frac{P(M_2|C_1)P(C_1)}{P(M_2)}
  =
  \frac{(1/2)(1/3)}{1/2}
  =
  \frac13.
\]
因此，给定 Monty 打开门 $2$，参赛者最初选择的门后有汽车的概率是 $1/3$，也就是说换门策略成功的概率是 $2/3$。

许多人第一次看到这个问题时，会认为换门没有优势：“剩下两扇门，其中一扇有汽车，所以概率是五五开。”上一章之后，我们知道这个论证误用了朴素概率定义。不过，即使朴素定义在这里不合适，它仍然强烈影响人的直觉。1990 年，Marilyn vos Savant 在《Parade》杂志专栏中给出 Monty Hall 问题的正确解法后，收到了成千上万封读者（甚至包括数学家）来信，坚持说她错了。

为了建立正确直觉，可以考虑一个极端情形。假设有一百万扇门，其中 $999999$ 扇门后是山羊，$1$ 扇门后是汽车。参赛者初选之后，Monty 打开 $999998$ 扇有山羊的门，并提供换门机会。在这个极端情形下，很清楚两扇未打开的门并不是五五开；很少有人会顽固地坚持最初选择。三门情形也是同样道理。

正如在例 2.2.6 中我们必须对如何遇到随机女孩作出假设一样，这里的换门策略成功率为 $2/3$，也依赖于我们对 Monty 如何决定打开哪扇门所作的假设。在习题中，我们会考虑 Monty Hall 问题的若干变体与推广，其中有些会改变换门策略是否值得采用。

\subsection*{2.7.2\quad 策略：以第一步为条件}

对于具有递归结构的问题，以实验的第一步为条件通常很有用。接下来两个例子都使用这一策略，我们称之为\textbf{第一步分析}。

\noindent\textbf{例 2.7.2（分枝过程）}\quad 一只名叫 Bobo 的变形虫生活在池塘中。一分钟后，Bobo 会等可能地死亡、分裂成两只变形虫，或保持不变；在随后的每一分钟里，所有存活的变形虫都会以同样方式独立行动。求变形虫种群最终灭绝的概率。

\begin{center}
  \includegraphics[width=0.42\textwidth]{2-5-1.png}
\end{center}

\noindent\textbf{解}\quad 令 $D$ 为种群最终灭绝的事件；我们要求 $P(D)$。以第一步结果为条件。令 $B_i$ 表示第一分钟后 Bobo 变成 $i$ 只变形虫的事件，其中 $i=0,1,2$。我们知道
\[
  P(D|B_0)=1,\qquad P(D|B_1)=P(D),
\]
因为如果 Bobo 保持不变，问题就回到原点。若 Bobo 分裂成两只，则我们有原问题的两个独立版本；要最终灭绝，两个后代都必须最终灭绝，所以
\[
  P(D|B_2)=P(D)^2.
\]
现在已经穷尽所有可能情形，可以用全概率公式合并：
\[
\begin{aligned}
  P(D)
  &=
  P(D|B_0)\cdot\frac13+
  P(D|B_1)\cdot\frac13+
  P(D|B_2)\cdot\frac13\\
  &=
  1\cdot\frac13+P(D)\cdot\frac13+P(D)^2\cdot\frac13.
\end{aligned}
\]
解得 $P(D)=1$：变形虫种群以概率 $1$ 灭绝。

第一步分析之所以在这里有效，是因为问题本质上是自相似的：当 Bobo 继续作为一只变形虫存在，或分裂成两只时，我们最终得到原问题的另一个版本，或两个版本。以第一步为条件，使我们能够用 $P(D)$ 本身来表示 $P(D)$。

\noindent\textbf{例 2.7.3（赌徒破产）}\quad 两个赌徒 A 与 B 进行一系列每局 $1$ 美元的赌博。每一局中，赌徒 A 以概率 $p$ 获胜，赌徒 B 以概率 $q=1-p$ 获胜。赌徒 A 初始有 $i$ 美元，赌徒 B 初始有 $N-i$ 美元；二人的总财富保持不变，因为每当 A 输掉 $1$ 美元，这 $1$ 美元就归 B，反之亦然。

我们可以把这个游戏看成整数 $0$ 到 $N$ 上的\textbf{随机游走}，其中 $p$ 是每一步向右走的概率：想象一个人从位置 $i$ 出发，每个时刻以概率 $p$ 向右走一步，以概率 $q=1-p$ 向左走一步。当 A 或 B 破产时，即随机游走到达 $0$ 或 $N$ 时，游戏结束。求 A 赢得游戏（带走所有钱）的概率。

\begin{center}
  \includegraphics[width=0.58\textwidth]{2-5-2.png}
\end{center}

\noindent\textbf{解}\quad 我们注意到，这个游戏像 Bobo 的繁殖过程一样具有递归结构：第一步之后，游戏与原来完全相同，只是 A 的财富变为 $i+1$ 或 $i-1$。令 $p_i$ 为在 A 初始有 $i$ 美元时 A 赢得游戏的概率。我们用第一步分析求解 $p_i$。令 $W$ 为 A 赢得游戏的事件。由全概率公式，以第一局结果为条件，有
\[
\begin{aligned}
  p_i
  &=
  P(W|\text{A 从 } i \text{ 开始，第一局胜})\cdot p
  +P(W|\text{A 从 } i \text{ 开始，第一局负})\cdot q\\
  &=
  P(W|\text{A 从 } i+1 \text{ 开始})\cdot p
  +P(W|\text{A 从 } i-1 \text{ 开始})\cdot q\\
  &=
  p_{i+1}p+p_{i-1}q.
\end{aligned}
\]
这对所有 $i=1,\ldots,N-1$ 都成立，并且边界条件为 $p_0=0,\ p_N=1$。现在可以解这个称为\textbf{差分方程}的方程来得到 $p_i$。数学附录 A.4 讨论了如何解差分方程；这里省略其中一些步骤。

该差分方程的特征方程为 $px^2-x+q=0$，其根为 $1$ 与 $q/p$。若 $p\ne 1/2$，这两个根不同，通解为
\[
  p_i=a\cdot 1^i+b\cdot\left(\frac{q}{p}\right)^i.
\]
利用边界条件 $p_0=0,\ p_N=1$，得
\[
  a=-b=\frac{1}{1-(q/p)^N}.
\]
代回即可得到特解。若 $p=1/2$，特征多项式的根不相异，因此通解为
\[
  p_i=a\cdot 1^i+b\cdot i\cdot 1^i.
\]
边界条件给出 $a=0,\ b=1/N$。

综上，A 以初始财富 $i$ 赢得游戏的概率为
\[
  p_i=
  \begin{cases}
    \dfrac{1-(q/p)^i}{1-(q/p)^N}, & p\ne 1/2,\\[1.2ex]
    \dfrac{i}{N}, & p=1/2.
  \end{cases}
\]
$p=1/2$ 的情形与 $p\ne 1/2$ 的情形相容，因为
\[
  \lim_{p\to 1/2}\frac{1-(q/p)^i}{1-(q/p)^N}=\frac{i}{N}.
\]
要看出这一点，令 $x=q/p$ 并令 $x$ 趋于 $1$。由 L'H\^opital 法则，
\[
  \lim_{x\to 1}\frac{1-x^i}{1-x^N}
  =
  \lim_{x\to 1}\frac{i x^{i-1}}{N x^{N-1}}
  =
  \frac{i}{N}.
\]

当 $p=1/2$ 时，答案有一个简单解释：A 获胜的概率等于 A 初始财富占总财富的比例。因此，如果 $p=1/2$ 且 A 初始财富远少于 B，那么 A 赢得游戏的机会很低。若 $p<1/2$，即使 $p$ 只比 $1/2$ 略小，且两人初始财富相同，A 获胜的机会也可能很低。例如，若 $p=0.49$，且两人各自以 $100$ 美元开始，则 A 获胜概率约为 $1.8\%$。

我们一直关注 A 赢得游戏的概率，那么 B 呢？不必从头开始，我们可以利用\textbf{对称性}：除记号外，游戏描述中没有任何东西区分 A 与 B。由对称性，B 以初始财富 $N-i$ 赢得游戏的概率，可以通过交换 $q$ 与 $p$、以及 $i$ 与 $N-i$ 的角色得到：
\[
  P(\text{B 获胜}|\text{B 从 }N-i\text{ 开始})
  =
  \begin{cases}
    \dfrac{1-(p/q)^{N-i}}{1-(p/q)^N}, & p\ne 1/2,\\[1.2ex]
    \dfrac{N-i}{N}, & p=1/2.
  \end{cases}
\]
于是可以验证，对所有 $i$ 与所有 $p$，有 $P(\text{A 获胜})+P(\text{B 获胜})=1$，所以游戏保证会结束：它永远振荡下去的概率为 $0$。

\section*{2.8\quad 陷阱与悖论}

接下来两个例子是在法律语境中出现的条件思考谬误。检察官谬误是把 $P(A|B)$ 与 $P(B|A)$ 混淆；辩护律师谬误则是没有以\textit{所有}证据为条件。

\noindent\textbf{生物危害 2.8.1（检察官谬误）}\quad 1998 年，Sally Clark 因两个儿子出生后不久死亡而被控谋杀。审判中，检方专家证人作证称，新生儿死于婴儿猝死综合征（SIDS）的概率为 $1/8500$，因此同一家庭中两名婴儿死于 SIDS 的概率为 $(1/8500)^2$，约为 $7300$ 万分之一。因此，他继续说，Clark 无罪的概率是 $7300$ 万分之一。

这一推理至少有两个重大问题。首先，专家证人通过相乘两个单独事件概率来求“第一个儿子死于 SIDS”与“第二个儿子死于 SIDS”的交集概率；正如我们所知，只有当同一家庭内 SIDS 死亡事件独立时，这样做才有效。如果遗传或其他家庭特定风险因素使某些家庭的新生儿都面临更高的 SIDS 风险，这种独立性就不成立。

其次，这位所谓专家混淆了两个不同的条件概率：$P(\text{无罪}|\text{证据})$ 不同于 $P(\text{证据}|\text{无罪})$。证人声称，如果被告无罪，观察到两名新生儿死亡的概率极低；也就是说，$P(\text{证据}|\text{无罪})$ 很小。但我们真正关心的是
\[
  P(\text{无罪}|\text{证据}),
\]
也就是给定所有证据时被告无罪的概率。由贝叶斯公式，
\[
  P(\text{无罪}|\text{证据})
  =
  \frac{
    P(\text{证据}|\text{无罪})P(\text{无罪})
  }{
    P(\text{证据})
  }.
\]
因此，为了计算给定证据后无罪的条件概率，必须把 $P(\text{无罪})$，即无罪的先验概率，考虑在内。这个概率极高：双重 SIDS 死亡很罕见，但双重杀婴同样罕见。把分母展开为
\[
  P(\text{证据}|\text{无罪})P(\text{无罪})
  +
  P(\text{证据}|\text{有罪})P(\text{有罪}),
\]
注意若 $P(\text{有罪})$ 足够小，使第二项相较第一项可以忽略，则 $P(\text{无罪}|\text{证据})$ 的分母近似等于分子，从而 $P(\text{无罪}|\text{证据})$ 接近 $1$。

给定证据后无罪的后验概率强烈依赖于 $P(\text{证据}|\text{无罪})$ 与 $P(\text{无罪})$ 这两个量，前者很低，后者很高。专家给出的 $(1/8500)^2$ 本身就值得怀疑，而且它只是整个方程的一部分。

遗憾的是，Clark 因谋杀罪被定罪并入狱，部分原因正是这位专家的错误证词；她在定罪被推翻前在狱中度过了三年多。Sally Clark 案中对条件概率的误用引发公愤，并导致数百起检方使用类似错误逻辑案件受到复查。

\noindent\textbf{生物危害 2.8.2（辩护律师谬误）}\quad 一名女子被谋杀，她的丈夫因这起犯罪受审。证据显示，被告过去曾虐待妻子。辩护律师主张，由于在虐待妻子的男子中，只有 $1/10000$ 后来会谋杀妻子，因此虐待证据应以不相关为由排除。法官是否应准许辩护律师的动议，将这项证据排除在审判之外？

假设辩护律师给出的 $1/10000$ 数字正确，并进一步假设，对于相关的夫妻总体，有以下概率：$1/10$ 的丈夫虐待妻子；$1/5$ 被谋杀的妻子是被丈夫谋杀的；谋杀妻子的丈夫中，$50\%$ 先前虐待过妻子。还假设，如果被谋杀妻子的丈夫\textit{没有}犯下谋杀罪，那么他虐待妻子的概率回到无条件虐待概率。

如何定义“相关总体”、以及如何估计这些概率，都是困难问题。例如，应当看城市、州、全国，还是国际统计？如何处理未报告的虐待和未侦破的谋杀？如果谋杀率随时间变化又如何？在本问题中，假设已经就相关总体的合理选择达成一致，并且所列概率已知为正确。

令 $A$ 为丈夫虐待妻子的事件，$G$ 为丈夫有罪的事件。辩护方的论点是 $P(G|A)=1/10000$，所以以既往虐待史为条件时，有罪仍极不可能。

然而，辩护律师没有以一个关键事实为条件：在本案中，我们\textit{知道妻子被谋杀了}。因此，相关概率不是 $P(G|A)$，而是 $P(G|A,M)$，其中 $M$ 为妻子被谋杀的事件。

带额外条件的贝叶斯公式给出
\[
\begin{aligned}
  P(G|A,M)
  &=
  \frac{P(A|G,M)P(G|M)}
       {P(A|G,M)P(G|M)+P(A|G^c,M)P(G^c|M)}\\
  &=
  \frac{0.5\cdot 0.2}{0.5\cdot 0.2+0.1\cdot 0.8}\\
  &=
  \frac59.
\end{aligned}
\]
所以有罪的后验概率 $P(G|A,M)$ 比辩护律师关注的 $P(G|A)$ 大 $5000$ 多倍。以虐待证据为条件，会把有罪概率从 $P(G|M)=0.2$ 提高到 $P(G|A,M)\approx 0.56$，因此被告的虐待史给出了非常重要的信息，恰与辩护律师的论点相反。

在上面对 $P(G|A,M)$ 的计算中，我们完全没有使用辩护律师的 $P(G|A)$ 数字；它与我们的计算无关，因为它没有考虑妻子已经被谋杀这一事实。我们必须以\textit{所有}证据为条件。

本章最后以一个关于条件概率与数据汇总的悖论收尾。

\noindent\textbf{例 2.8.3（Simpson 悖论）}\quad Hibbert 医生与 Nick 医生各自进行两类手术：心脏手术和创可贴移除。每次手术可能成功，也可能失败。两位医生的记录如下表，并在图 2.6 中以图形方式展示，其中白点表示成功手术，黑点表示失败手术。

\[
\begin{array}{c|cc}
 & \text{心脏} & \text{创可贴}\\
\hline
\text{成功} & 70 & 10\\
\text{失败} & 20 & 0
\end{array}
\qquad
\begin{array}{c|cc}
 & \text{心脏} & \text{创可贴}\\
\hline
\text{成功} & 2 & 81\\
\text{失败} & 8 & 9
\end{array}
\]
\[
  \text{Hibbert 医生}
  \qquad\qquad\qquad
  \text{Nick 医生}
\]

Hibbert 医生在心脏手术中的成功率高于 Nick 医生：$90$ 次中成功 $70$ 次，相比于 $10$ 次中成功 $2$ 次。Hibbert 医生在创可贴移除中的成功率也更高：$10$ 次中成功 $10$ 次，相比于 $90$ 次中成功 $81$ 次。但若跨两类手术汇总来比较总体手术成功率，Hibbert 医生在 $100$ 次手术中成功 $80$ 次，而 Nick 医生在 $100$ 次手术中成功 $83$ 次：Nick 医生的总体成功率更高。

\begin{figure}[htbp]
  \centering
  \includegraphics[width=0.62\textwidth]{2-6.png}
  \caption{Simpson 悖论的一个例子。白点表示成功手术，黑点表示失败手术。Hibbert 医生在两类手术中都更好，但总体成功率更低，因为他进行更困难手术的频率远高于 Nick 医生。}
\end{figure}

这里发生的事情是，Hibbert 医生大概因为作为更优秀医生的声誉，进行了更多心脏手术；而心脏手术本来就比创可贴移除风险更高。他的总体成功率较低，不是因为他在任何一种特定手术上的技术较差，而是因为他所做手术中风险较高的比例更大。

用事件记号可以把这一点说清楚。对事件 $A,B,C$，若
\[
  P(A|B,C)<P(A|B^c,C),
\]
\[
  P(A|B,C^c)<P(A|B^c,C^c),
\]
但
\[
  P(A|B)>P(A|B^c),
\]
则称出现了 Simpson 悖论。

在本例中，令 $A$ 为手术成功的事件，$B$ 为 Nick 医生是外科医生的事件，$C$ 为手术是心脏手术的事件。Simpson 悖论的条件成立，因为无论以心脏手术还是创可贴移除为条件，Nick 医生的手术成功概率都低于 Hibbert 医生；但总体成功概率却是 Nick 医生更高。

全概率公式说明了为什么这种现象会在数学上发生：
\[
\begin{aligned}
  P(A|B)&=P(A|C,B)P(C|B)+P(A|C^c,B)P(C^c|B),\\
  P(A|B^c)&=P(A|C,B^c)P(C|B^c)+P(A|C^c,B^c)P(C^c|B^c).
\end{aligned}
\]
上面的方程把 $P(A|B)$ 表示为 $P(A|C,B)$ 与 $P(A|C^c,B)$ 的加权平均，把 $P(A|B^c)$ 表示为 $P(A|C,B^c)$ 与 $P(A|C^c,B^c)$ 的加权平均。如果两个加权平均中的相应权重相同，Simpson 悖论就不可能发生。但这里的权重不同：
\[
  P(C|B)<P(C|B^c),\qquad P(C^c|B)>P(C^c|B^c),
\]
因为 Nick 医生做心脏手术的概率远低于 Hibbert 医生。

虽然我们有
\[
  P(A|C,B)<P(A|C,B^c)
\]
以及
\[
  P(A|C^c,B)<P(A|C^c,B^c),
\]
但权重差异如此之大，使得在不以 $C$ 是否发生为条件时，不等式方向翻转：
\[
  P(A|B)>P(A|B^c).
\]
数值上，两个加权平均为
\[
  P(A|B)=0.83=(2/10)\cdot 0.1+(81/90)\cdot 0.9,
\]
\[
  P(A|B^c)=0.80=(70/90)\cdot 0.9+(10/10)\cdot 0.1.
\]
第一个方程（对应 Nick 医生）在第二项（对应较容易手术）上赋予了比第二个方程更大的权重。

跨不同手术类型汇总，会给医生能力带来误导性图像，因为我们丢失了哪位医生倾向于做哪类手术的信息。当我们认为手术类型这样的\textbf{混杂变量}可能起作用时，应当检查分组后的数据，看看真正发生了什么。

Simpson 悖论出现在许多现实语境中。下面这些例子中，你应当尝试识别出产生悖论的事件 $A,B,C$。

\begin{itemize}
  \item \textit{大学录取中的性别歧视：}20 世纪 70 年代，在加州大学伯克利分校研究生录取中，男性被录取的可能性显著高于女性，因此引发性别歧视指控。然而，在大多数单个院系内，女性录取率高于男性。后来发现，女性更倾向于申请录取竞争更激烈的院系，而男性更倾向于申请竞争较不激烈的院系。
  \item \textit{棒球打击率：}球员 1 可能在棒球赛季前半程的打击率高于球员 2，并且在后半程打击率也高于球员 2，但整个赛季的总体打击率反而更低。这取决于两名球员在赛季两个半程中各有多少次击球机会。所谓一次击球，是球员轮到尝试击球；球员的打击率等于安打数除以击球次数。
  \item \textit{吸烟对健康的影响：}Cochran 的研究发现，在任意年龄组内，吸烟者的死亡率都高于抽雪茄者；但由于吸烟者平均比抽雪茄者年轻，所以吸烟者的总体死亡率反而更低。
\end{itemize}

\section*{2.9\quad 回顾}

给定 $B$ 时 $A$ 的条件概率为
\[
  P(A|B)=\frac{P(A\cap B)}{P(B)}.
\]

条件概率具有与概率完全相同的性质，但 $P(\cdot|B)$ 会更新我们对事件的不确定性，以反映已经观察到的证据 $B$。若某事件在观察到证据 $B$ 后概率不变，则称它与 $B$ 独立。两个事件还可以在给定第三个事件 $E$ 时条件独立。条件独立不推出无条件独立，无条件独立也不推出条件独立。

关于条件概率的两个重要结果是贝叶斯公式和全概率公式。贝叶斯公式联系 $P(A|B)$ 与 $P(B|A)$，全概率公式则允许我们把样本空间划分成若干部分，在每一部分中计算条件概率，再合并得到无条件概率。

贝叶斯公式说
\[
  P(A|B)=\frac{P(B|A)P(A)}{P(B)},
\]
而全概率公式说
\[
  P(B)=\sum_{i=1}^n P(B|A_i)P(A_i),
\]
其中 $A_1,\ldots,A_n$ 是样本空间的任一划分。贝叶斯公式与全概率公式常常配合使用。

条件化对于解题极其有用，因为它让我们可以把一个问题分解为更小的片段，分别考虑所有可能情形，然后再合并。当使用这种策略时，应尝试以“若知道便会使问题更简单”的信息为条件，因此有“以你希望知道的信息为条件”的说法。当一个问题包含多个阶段时，以第一步为条件通常有助于得到递归关系。

条件思考中的常见错误包括：
\begin{itemize}
  \item 混淆先验概率 $P(A)$ 与后验概率 $P(A|B)$；
  \item 检察官谬误，即混淆 $P(A|B)$ 与 $P(B|A)$；
  \item 辩护律师谬误，即未以所有证据为条件；
  \item 忽视 Simpson 悖论，以及没有仔细思考是否应汇总数据。
\end{itemize}

图 2.7 说明了当新证据依次到来时，概率如何更新。设想我们关心某个事件 $A$。例如，周一早晨我们对 $A$ 的先验概率是 $P(A)$。如果周一下午观察到 $B$ 发生，那么可以用贝叶斯公式（或条件概率定义）计算后验概率 $P(A|B)$。

我们把这个关于 $A$ 的后验概率作为周二早晨的新先验概率，然后继续收集证据。假设周二观察到 $C$ 发生，那么可以用多种方式计算新的后验概率 $P(A|B,C)$（在这个语境下，最自然的方式大概是使用带额外条件 $B$ 的贝叶斯公式）。如果还要继续收集证据，这个新的后验概率又会成为新的先验概率。

\begin{figure}[htbp]
  \centering
  \includegraphics[width=0.72\textwidth]{2-7.png}
  \caption{条件概率说明了如何在新证据到来时更新概率。图中展示了事件 $A$ 最初的概率、获得一条证据 $B$ 后的概率，以及获得第二条证据 $C$ 后的概率。观察到第一条证据后的后验概率，会在观察第二条证据前成为新的先验概率；在 $B$ 与 $C$ 都被观察到后，可以用多种方式求 $A$ 的新后验概率。若继续收集更多证据，它又会成为新的先验概率。}
\end{figure}

\section*{2.11\quad 习题}

标有（S）的习题在 \texttt{http://stat110.net} 有详细解答。

\subsection*{以证据为条件}

\begin{enumerate}
  \item （S）一个垃圾邮件过滤器通过考察垃圾邮件中常见的短语来设计。假设 $80\%$ 的电子邮件是垃圾邮件。在 $10\%$ 的垃圾邮件中会使用短语“free money”，而在非垃圾邮件中只有 $1\%$ 会使用该短语。现在收到一封新邮件，其中确实提到“free money”。求这封邮件是垃圾邮件的概率。

  \item （S）一名女子怀了双胞胎男孩。双胞胎可能是同卵，也可能是异卵。假设出生的双胞胎中有 $1/3$ 是同卵；同卵双胞胎有 $50\%$ 的概率都是男孩，$50\%$ 的概率都是女孩；异卵双胞胎中，每个孩子独立地有 $50\%$ 的概率是男孩，$50\%$ 的概率是女孩。给定上述信息，求这名女子的双胞胎是同卵的概率。

  \item 根据 CDC（美国疾病控制与预防中心）的资料，吸烟男性患肺癌的可能性是不吸烟男性的 $23$ 倍。CDC 还称，美国男性中有 $21.6\%$ 吸烟。给定一名美国男性患肺癌，求他是吸烟者的概率。

  \item Fred 正在回答考试中的一道选择题，需要在 $n$ 个选项中选择一个（其中恰有一个正确）。令 $K$ 为 Fred 知道答案的事件，$R$ 为他答对题目的事件（无论是凭知识还是运气）。假设如果他知道正确答案，就一定答对；如果不知道，就完全随机猜测。设 $P(K)=p$。
  \begin{enumerate}
    \item 求 $P(K|R)$（用 $p,n$ 表示）。
    \item 证明 $P(K|R)\ge p$，并解释这为什么直观上合理。何时（如果有的话）$P(K|R)=p$？
  \end{enumerate}

  \item 从一副充分洗匀的标准扑克牌中发三张牌。前两张被翻开，显示第一张是黑桃 A，第二张是梅花 8。给定这些信息，分别用条件概率定义和对称性求第三张牌是 A 的概率。

  \item 一个帽子里有 $100$ 枚硬币，其中 $99$ 枚公平，$1$ 枚是双正面硬币。随机均匀选一枚硬币，将它抛掷 $7$ 次，结果 $7$ 次全为正面。给定这一信息，求所选硬币是双正面硬币的概率。

  \item 一个帽子里有 $100$ 枚硬币，其中至少 $99$ 枚公平，但也可能有一枚双正面硬币；若没有这枚双正面硬币，则帽子中有 $100$ 枚公平硬币。令 $D$ 为存在这样一枚双正面硬币的事件，并假设 $P(D)=1/2$。随机均匀选一枚硬币，将它抛掷 $7$ 次，结果 $7$ 次全为正面。
  \begin{enumerate}
    \item 给定这一信息，存在双正面硬币的概率是多少？
    \item 给定这一信息，所选硬币是双正面硬币的概率是多少？
  \end{enumerate}

  \item 某类手机屏幕由三家公司 A、B、C 制造。三家公司供应屏幕的比例分别为 $0.5,0.3,0.2$；其屏幕有缺陷的概率分别为 $0.01,0.02,0.03$。给定某手机的屏幕有缺陷，求它由公司 A 制造的概率。

  \item
  \begin{enumerate}
    \item 证明：若事件 $A_1$ 与 $A_2$ 有相同先验概率 $P(A_1)=P(A_2)$，且 $A_1$ 蕴含 $B$、$A_2$ 也蕴含 $B$，那么观察到 $B$ 发生后，二者有相同后验概率：
    \[
      P(A_1|B)=P(A_2|B).
    \]
    \item 解释为什么 (a) 直观上合理，并给出一个具体例子。
  \end{enumerate}

  \item Fred 正在做一个大型项目。项目计划中设置了两个里程碑，并规定了它们应完成的日期，以便追踪 Fred 的进度。令 $A_1$ 为 Fred 按时完成第一个里程碑的事件，$A_2$ 为他按时完成第二个里程碑的事件，$A_3$ 为他按时完成项目的事件。

  假设对 $j=1,2$ 有
  \[
    P(A_{j+1}|A_j)=0.8,\qquad P(A_{j+1}|A_j^c)=0.3,
  \]
  因为如果 Fred 落后于进度，就很难赶上。另外，假设第二个里程碑取代第一个：一旦知道 Fred 是否按时完成第二个里程碑，第一个里程碑发生了什么就不再重要。可表达为：$A_1$ 与 $A_3$ 在给定 $A_2$ 时条件独立，并且在给定 $A_2^c$ 时也条件独立。
  \begin{enumerate}
    \item 求给定 Fred 按时完成第一个里程碑时，他按时完成项目的概率；再求给定 Fred 未按时完成第一个里程碑时，他按时完成项目的概率。
    \item 假设 $P(A_1)=0.75$。求 Fred 按时完成项目的概率。
  \end{enumerate}

  \item 选举出口民调是在选民刚投完票后进行的调查。出口民调的一个主要用途，是让新闻机构能在正式计票完成前尽快判断谁赢得选举。这在若干选举中出过严重错误，有时是因为\textbf{选择偏差}：被邀请并同意参加调查的人样本，可能与总体选民不够相似。

  考虑一场有候选人 A 与候选人 B 的选举。每位选民都被邀请参加出口民调，并被询问投给谁；有些人接受，有些人拒绝。随机选一名选民，令 $A$ 为其投给 A 的事件，$W$ 为其愿意参加出口民调的事件。假设 $P(W|A)=0.7$，但 $P(W|A^c)=0.3$。在出口民调中，$60\%$ 的受访者说他们投给了 A（假设他们都诚实），这看起来表示 A 会轻松获胜。求 $P(A)$，也就是实际投给 A 的人群比例。

  \item Alice 试图通过一个信道向 Bob 发送用二进制编码的信息。
  \begin{enumerate}
    \item 本小题中，假设她只发送一位（0 或 1），且两者概率相等。如果她发送 0，有 $5\%$ 的概率发生错误，导致 Bob 收到 1；如果她发送 1，有 $10\%$ 的概率发生错误，导致 Bob 收到 0。给定 Bob 收到 1，求 Alice 实际发送 1 的概率。
    \item 为降低误通信概率，Alice 与 Bob 决定使用\textbf{重复码}。Alice 仍想传达 0 或 1，但这次她额外重复两次，即发送 000 表示 0，发送 111 表示 1。Bob 将按多数位解码。假设错误概率与 (a) 相同，且不同位上的错误事件相互独立。给定 Bob 收到 110，求 Alice 原本想传达 1 的概率。
  \end{enumerate}

  \item 公司 A 刚开发出一种针对某疾病的诊断测试。该疾病影响 $1\%$ 的人群。如例 2.3.9 中定义，测试的灵敏度是有病者检测为阳性的概率，特异度是无病者检测为阴性的概率。假设与例 2.3.9 一样，灵敏度和特异度均为 $0.95$。

  公司 B 是公司 A 的竞争者，也提供该疾病的竞争性检测。公司 B 声称其检测比公司 A 更快、更便宜、更不痛苦（公司 A 的检测需要切口），而且有更高的总体成功率；这里总体成功率定义为随机一人被正确诊断的概率。
  \begin{enumerate}
    \item 事实证明，公司 B 的检测可以极其简单地描述和执行：无论病人是谁，都诊断为没有该病。检验公司 B 关于总体成功率的主张是否为真。
    \item 解释为什么公司 A 的检测仍然可能有用。
    \item 公司 A 想开发一种新检测，使其总体成功率高于公司 B 的检测。若灵敏度与特异度相等，那么灵敏度必须多高才能达到目标？如果（令人惊讶地）能使灵敏度等于 $1$，那么特异度必须多高？如果（令人惊讶地）能使特异度等于 $1$，那么灵敏度必须多高？
  \end{enumerate}

  \item 考虑 Tversky 与 Kahneman 给出的情境：令 $A$ 为 Peter 在明年年底前在家中安装防盗报警系统的事件，令 $B$ 为 Peter 家在明年年底前被盗的事件。
  \begin{enumerate}
    \item 直观上，你认为 $P(A|B)$ 与 $P(A|B^c)$ 哪个更大？解释你的直觉。
    \item 直观上，你认为 $P(B|A)$ 与 $P(B|A^c)$ 哪个更大？解释你的直觉。
    \item 证明：对任意事件 $A,B$（概率都不等于 $0$ 或 $1$），不等式 $P(A|B)>P(A|B^c)$ 等价于 $P(B|A)>P(B|A^c)$。
    \item Tversky 与 Kahneman 报告说，在他们提出 (a) 与 (b) 的 $162$ 人中，有 $131$ 人认为 $P(A|B)>P(A|B^c)$ 且 $P(B|A)<P(B|A^c)$。尽管 (c) 说明这两个不等式不可能同时成立，为什么这种意见会如此流行？给出一个合理解释。
  \end{enumerate}

  \item 设 $A,B$ 为事件，且
  \[
    0<P(A\cap B)<P(A)<P(B)<P(A\cup B)<1.
  \]
  你希望 $A$ 与 $B$ 都发生。下列信息中，你最希望观察到哪一个：$A$ 发生、$B$ 发生，还是 $A\cup B$ 发生？

  \item 证明：$P(A|B)\le P(A)$ 蕴含 $P(A|B^c)\ge P(A)$，并给出直观解释。

  \item 在确定性逻辑中，“$A$ 蕴含 $B$”等价于它的逆否命题“非 $B$ 蕴含非 $A$”。本题考虑概率中的类似陈述。令 $A,B$ 为概率都不等于 $0$ 或 $1$ 的事件。
  \begin{enumerate}
    \item 证明：若 $P(B|A)=1$，则 $P(A^c|B^c)=1$。

    提示：使用贝叶斯公式与全概率公式。
    \item 但证明：若把 (a) 中的等号换成近似等号，则结论一般不成立。特别地，给出一个例子，使 $P(B|A)$ 非常接近 $1$，但 $P(A^c|B^c)$ 非常接近 $0$。

    提示：若 $A$ 与 $B$ 独立，会发生什么？
  \end{enumerate}

  \item 证明：若 $P(A)=1$，则对任何满足 $P(B)>0$ 的 $B$，有 $P(A|B)=1$。直观上，这表示若某人武断地以绝对确定性相信某事，那么再多证据都不会改变其看法。避免把概率 $0$ 或 $1$ 赋给任何事件（数学确定性除外）的原则被统计学家 Dennis Lindley 称为 Cromwell 法则，因为 Cromwell 曾对苏格兰教会说：“想想你们也可能错了。”

  提示：写出 $P(B)=P(B\cap A)+P(B\cap A^c)$，然后证明 $P(B\cap A^c)=0$。

  \item 用条件概率解释 Sherlock Holmes 的名言，并仔细区分先验概率与后验概率：“这是我的一条老格言：当你排除不可能之后，剩下的无论多么不可能，都必然是真相。”

  \item 从一副牌中取出黑桃 J（带苹果酒）、红桃 J（带馅饼）、黑桃 Q（带眨眼）和红桃 Q（没有馅饼）。洗匀这四张牌，然后发两张。注意：解题时不需要考虑苹果酒、馅饼和眨眼的字面含义。
  \begin{enumerate}
    \item 给定第一张牌是 Q，求两张牌都是 Q 的概率。
    \item 给定至少一张牌是 Q，求两张牌都是 Q 的概率。
    \item 给定其中一张牌是红桃 Q，求两张牌都是 Q 的概率。
  \end{enumerate}

  \item 将一枚公平硬币抛掷 $3$ 次。每次抛掷结果分别记录在一张纸条上（正面写 H，反面写 T），然后把 $3$ 张纸条放入帽子中。
  \begin{enumerate}
    \item 给定至少 $2$ 次为正面，求 $3$ 次全为正面的概率。
    \item 从帽子中随机抽出两张纸条，二者都显示 H。给定这一信息，求 $3$ 次全为正面的概率。
  \end{enumerate}

  \item （S）一个袋子里有一颗弹珠，它等可能地是绿色或蓝色。向袋中放入一颗绿色弹珠（因此现在有两颗弹珠），然后随机取出一颗。取出的弹珠是绿色。求剩余弹珠也是绿色的概率。

  \item （S）令 $G$ 为某人犯下某起抢劫案的事件。在搜集证据时，得知事件 $E_1$ 发生；稍后又得知另一个事件 $E_2$ 发生。是否可能出现这样的情形：单独看时，每条证据都会增加有罪概率，即 $P(G|E_1)>P(G)$ 且 $P(G|E_2)>P(G)$，但合在一起却降低有罪概率，即 $P(G|E_1,E_2)<P(G)$？

  \item 是否可能存在事件 $A_1,A_2,B,C$，使得 $P(A_1|B)>P(A_1|C)$ 且 $P(A_2|B)>P(A_2|C)$，但 $P(A_1\cup A_2|B)<P(A_1\cup A_2|C)$？若可能，给出一个例子（用故事解释事件，并给出具体数字）；否则，证明这种现象不可能发生。

  \item （S）一场犯罪由两个嫌疑人 $A$ 与 $B$ 中的一人实施。起初，对二人的证据相同。进一步调查犯罪现场时发现，罪犯具有一种在 $10\%$ 人群中出现的血型。嫌疑人 $A$ 符合这种血型，而嫌疑人 $B$ 的血型未知。
  \begin{enumerate}
    \item 给定这条新信息，$A$ 是罪犯的概率是多少？
    \item 给定这条新信息，$B$ 的血型也与犯罪现场发现的血型相同的概率是多少？
  \end{enumerate}

  \item （S）为对抗垃圾邮件，Bob 安装了两个反垃圾邮件程序。一封邮件到达，它要么是正常邮件（事件 $L$），要么是垃圾邮件（事件 $L^c$）；程序 $j$ 要么把它标为正常（事件 $M_j$），要么标为垃圾（事件 $M_j^c$），其中 $j\in\{1,2\}$。假设 Bob 的邮件中有 $10\%$ 是正常邮件，并且两个程序在如下意义上都“$90\%$ 准确”：
  \[
    P(M_j|L)=P(M_j^c|L^c)=\frac9{10}.
  \]
  还假设给定一封邮件是否为垃圾，两程序输出条件独立。
  \begin{enumerate}
    \item 给定第一个程序将邮件标为正常，求该邮件为正常邮件的概率（化简）。
    \item 给定两个程序都将邮件标为正常，求该邮件为正常邮件的概率（化简）。
    \item Bob 运行第一个程序且 $M_1$ 发生。他更新概率，然后运行第二个程序。令 $\widetilde P(A)=P(A|M_1)$ 为运行第一个程序后的更新概率函数。简要用文字说明 $\widetilde P(L|M_2)=P(L|M_1\cap M_2)$ 是否成立：一次性以 $M_1\cap M_2$ 为条件，是否等价于先以 $M_1$ 为条件更新概率，再以 $M_2$ 为条件？
  \end{enumerate}

  \item 假设总体中有 $5$ 种血型，分别记为类型 $1$ 到 $5$，概率为 $p_1,p_2,\ldots,p_5$。一场犯罪由两个人实施。某嫌疑人血型为类型 $1$，其先验有罪概率为 $p$。在犯罪现场收集到的血液证据显示，两名罪犯中一人血型为类型 $1$，另一人血型为类型 $2$。

  给定证据，求该嫌疑人有罪的后验概率。证据是否使其更可能有罪、更不可能有罪，还是取决于参数 $p,p_1,\ldots,p_5$ 的取值？若取决于这些值，请给出证据使其更可能有罪的简单判别准则。

  \item Fred 刚刚对某疾病检测为阳性。
  \begin{enumerate}
    \item 给定这一信息，用先验赔率、测试灵敏度与测试特异度表示他患病的后验赔率。
    \item Fred 自然更关心 $P(\text{患病}|\text{检测阳性})$，也就是阳性预测值，而不是灵敏度 $P(\text{检测阳性}|\text{患病})$。生物统计学与流行病学中有一条方便的经验法则：

    \textit{对于罕见疾病和相当好的检测，特异度在决定阳性预测值时比灵敏度重要得多。}

    直观解释这条经验法则为什么有效。本小题中你可以自行设定一些具体数字，并用频率学派方式把概率解释为大人群中的比例，例如假设一种疾病影响 $10000$ 人中的 $1\%$，然后考虑灵敏度与特异度的各种可能值。
  \end{enumerate}

  \item 一个家庭有两个孩子。令 $C$ 为孩子可能具有的某种特征，并假设每个孩子以概率 $p$ 具有特征 $C$，且彼此独立，也与性别独立。例如，$C$ 可以是例 2.2.7 中的“冬季出生”。在例 2.2.5 的假设下，证明：给定至少一个孩子是具有特征 $C$ 的女孩，两个孩子都是女孩的概率为
  \[
    \frac{2-p}{4-p}.
  \]
  注意当 $p=1$ 时这等于 $1/3$（与例 2.2.5 第一部分一致），而当 $p\to 0$ 时它从下方趋于 $1/2$（与例 2.2.7 一致）。

\subsection*{独立性与条件独立性}

  \item （S）一个家庭有 $3$ 个孩子，创造性地命名为 $A,B,C$。
  \begin{enumerate}
    \item 直观而清楚地讨论事件“$A$ 比 $B$ 年长”是否独立于事件“$A$ 比 $C$ 年长”。
    \item 给定 $A$ 比 $C$ 年长，求 $A$ 比 $B$ 年长的概率。
  \end{enumerate}

  \item （S）一个事件是否可能与自身独立？若可能，何时如此？

  \item （S）考虑四个非标准骰子（Efron 骰子），其面标如下（每个骰子的 $6$ 个面等可能）：
  \[
  \begin{array}{ll}
    A: & 4,4,4,4,0,0,\\
    B: & 3,3,3,3,3,3,\\
    C: & 6,6,2,2,2,2.
  \end{array}
  \]
  \[
  \begin{array}{ll}
    D: & 5,5,5,1,1,1.
  \end{array}
  \]
  四个骰子各掷一次。令 $A$ 为骰子 $A$ 的结果，$B$ 为骰子 $B$ 的结果，依此类推。
  \begin{enumerate}
    \item 求 $P(A>B),P(B>C),P(C>D)$ 和 $P(D>A)$。
    \item 事件 $A>B$ 是否独立于事件 $B>C$？事件 $B>C$ 是否独立于事件 $C>D$？解释。
  \end{enumerate}

  \item Alice、Bob 和另外 $100$ 个人住在一个小镇上。令 $C$ 为这另外 $100$ 个人组成的集合，令 $A$ 为 $C$ 中与 Alice 是朋友的人组成的集合，令 $B$ 为 $C$ 中与 Bob 是朋友的人组成的集合。假设对 $C$ 中每个人，Alice 以概率 $1/2$ 与其为朋友；Bob 同理；所有这些友谊状态相互独立。
  \begin{enumerate}
    \item 令 $D\subseteq C$。求 $P(A=D)$。
    \item 求 $P(A\subseteq B)$。
    \item 求 $P(A\cup B=C)$。
  \end{enumerate}

  \item 假设司机分为好司机和坏司机两类。令 $G$ 为某男子是好司机的事件，$A$ 为他明年发生车祸的事件，$B$ 为他后年发生车祸的事件。设
  \[
    P(G)=g,\quad P(A|G)=P(B|G)=p_1,\quad
    P(A|G^c)=P(B|G^c)=p_2,
  \]
  且 $p_1<p_2$。假设给定该男子是否为好司机后，$A$ 与 $B$ 独立（为简单起见，也为了避免过于沉重，假设这些车祸都是小事故，不会使他不能开车）。
  \begin{enumerate}
    \item 直观解释 $A$ 与 $B$ 是否独立。
    \item 求 $P(G|A^c)$。
    \item 求 $P(B|A^c)$。
  \end{enumerate}

  \item （S）你将和一位此前从未交手的对手下两盘棋。就本题而言，你的对手等可能是初学者、中级棋手或大师。相应地，你赢得单盘的概率分别为 $90\%,50\%,30\%$。
  \begin{enumerate}
    \item 你赢得第一盘的概率是多少？
    \item 恭喜你赢了第一盘！给定这一信息，你赢得第二盘的概率是多少？（假设给定对手水平后，两盘棋结果独立。）
    \item 解释“假设两盘棋结果独立”与“假设给定对手水平后两盘棋结果条件独立”的区别。哪一个假设更合理？为什么？
  \end{enumerate}

  \item
  \begin{enumerate}
    \item 假设在大学申请者总体中，棒球打得好与某个标准化考试数学成绩好相互独立（相对于某种“好”的度量）。某大学采用简单录取规则：当且仅当申请者棒球好或数学成绩好时录取。

    直观解释为什么在该大学录取的学生中，数学成绩好与棒球好负相关是合理的；也就是说，在已被录取的条件下，数学成绩好会降低棒球好的机会。

    \item 证明：若 $A$ 与 $B$ 独立且 $C=A\cup B$，则给定 $C$ 后 $A$ 与 $B$ 条件依赖（只要 $P(A\cap B)>0$ 且 $P(A\cup B)<1$），且
    \[
      P(A|B,C)<P(A|C).
    \]
    这种现象称为 Berkson 悖论，尤其常见于学校、医院等录取或接纳语境。
  \end{enumerate}

  \item 两种不同疾病会导致某种奇怪症状；任何患有其中一种或两种疾病的人都会出现该症状。令 $D_1$ 为患第一种疾病的事件，$D_2$ 为患第二种疾病的事件，$W$ 为出现奇怪症状的事件。假设 $D_1$ 与 $D_2$ 独立，且 $P(D_j)=p_j$；一个两种病都没有的人以概率 $w_0$ 出现该症状。令 $q_j=1-p_j$，并假设 $0<p_j<1$。
  \begin{enumerate}
    \item 求 $P(W)$。
    \item 求 $P(D_1|W),P(D_2|W)$ 和 $P(D_1,D_2|W)$。
    \item 用代数判断 $D_1$ 与 $D_2$ 在给定 $W$ 时是否条件独立。
    \item 仅本小题假设 $w_0=0$。用文字给出清楚且有说服力的直观解释，说明 $D_1$ 与 $D_2$ 在给定 $W$ 时是否条件独立。
  \end{enumerate}

  \item 我们想为电子邮件设计垃圾邮件过滤器。如习题 1 所述，一个主要策略是寻找在垃圾邮件中比在非垃圾邮件中更可能出现的短语。该习题中只考虑一个短语“free money”。更现实地，假设我们已经创建了一个由 $100$ 个词或短语组成的列表，它们在垃圾邮件中比在非垃圾邮件中更可能出现。

  令 $W_j$ 为电子邮件包含列表中第 $j$ 个词或短语的事件。令
  \[
    p=P(\text{垃圾邮件}),\quad
    p_j=P(W_j|\text{垃圾邮件}),\quad
    r_j=P(W_j|\text{非垃圾邮件}),
  \]
  其中“垃圾邮件”是该邮件为垃圾邮件事件的简写。

  假设 $W_1,\ldots,W_{100}$ 在给定邮件为垃圾邮件时条件独立，并且在给定邮件不是垃圾邮件时也条件独立。基于这类假设对电子邮件（或其他对象）分类的方法称为\textbf{朴素贝叶斯分类器}。（这里“朴素”指条件独立性是很强的假设，不是说贝叶斯本人朴素。这个假设可能现实，也可能不现实，但朴素贝叶斯分类器有时在实践中即使假设不现实也能工作得很好。）

  在这个假设下，例如我们知道
  \[
    P(W_1,W_2,W_3^c,W_4^c,\ldots,W_{100}^c|\text{垃圾邮件})
    =
    p_1p_2(1-p_3)(1-p_4)\cdots(1-p_{100}).
  \]
  如果没有朴素贝叶斯假设，就会有远为巨大的统计与计算困难，因为我们需要考虑 $2^{100}\approx 1.3\times 10^{30}$ 个形如
  \[
    A_1\cap A_2\cap\cdots\cap A_{100}
  \]
  的事件，其中每个 $A_j$ 等于 $W_j$ 或 $W_j^c$。现在一封新邮件到达，它包含列表中第 $23,64,65$ 个词或短语（但不包含其他 $97$ 个）。因此我们想计算
  \[
  P(\text{垃圾邮件}|W_1^c,\ldots,W_{22}^c,W_{23},
  W_{24}^c,\ldots,W_{63}^c,W_{64},W_{65},W_{66}^c,\ldots,W_{100}^c).
  \]
  注意，我们需要以\textit{所有}证据为条件，而不仅仅是 $W_{23}\cap W_{64}\cap W_{65}$ 已发生这一事实。求这封新邮件是垃圾邮件的条件概率（用 $p,p_j,r_j$ 表示）。

\subsection*{Monty Hall}

  \item （S）
  \begin{enumerate}
    \item 考虑 Monty Hall 问题的 $7$ 门版本。有 $7$ 扇门，其中一扇门后有汽车（你想要），其余门后有山羊（你不想要）。起初，汽车所在位置的所有可能性等可能。你选择一扇门。随后 Monty Hall 打开 $3$ 扇有山羊的门，并让你选择是否换到剩下 $3$ 扇门中的任意一扇。

    假设 Monty Hall 知道哪扇门后有汽车，总会打开 $3$ 扇山羊门并提供换门机会，而且在所有可打开的山羊门选择中等可能地选择。你应该换门吗？如果换到剩下 $3$ 扇门中的一扇，成功概率是多少？
    \item 将上题推广到 $n\ge 3$ 扇门的 Monty Hall 问题，其中 Monty 打开 $m$ 扇山羊门，$1\le m\le n-2$。
  \end{enumerate}

  \item （S）考虑 Monty Hall 问题的一个变体：Monty 比起打开门 $3$ 更喜欢打开门 $2$；如果他可以在打开这两扇门之间选择，就以概率 $p$ 打开门 $2$，其中 $1/2\le p\le 1$。

  回顾：有三扇门，其中一扇门后有汽车（你想要），另外两扇门后有山羊（你不想要）。起初，汽车所在位置三种可能性等可能。你选择一扇门，为具体起见假设是门 $1$。Monty 随后打开一扇门展示山羊，并提供换门机会。假设 Monty Hall 知道哪扇门后有汽车，总会打开一扇山羊门并提供换门机会；并且如上所述，若他可在打开门 $2$ 与门 $3$ 之间选择，他以概率 $p$ 打开门 $2$（其中 $1/2\le p\le 1$）。
  \begin{enumerate}
    \item 求总是换门策略成功的无条件概率（“无条件”是指不以 Monty 打开门 $2$ 还是门 $3$ 为条件）。
    \item 给定 Monty 打开门 $2$，求总是换门策略成功的概率。
    \item 给定 Monty 打开门 $3$，求总是换门策略成功的概率。
  \end{enumerate}

  \item Monty Hall 节目的收视率略有下降，一位惊慌的执行制片人向 Monty 抱怨，说节目中他开门的部分缺乏悬念：Monty 总是打开一扇有山羊的门。Monty 回答说这样做是为了不因展示汽车而毁掉游戏，但他同意如下更新游戏。

  每场节目之前，Monty 秘密抛一枚正面概率为 $p$ 的硬币。如果硬币为正面，Monty 决定打开一扇有山羊的门（若有选择则等可能选择）。否则，Monty 决定等可能随机打开一扇门。当然，Monty 不会打开参赛者最初选择的门。参赛者知道 $p$，但不知道抛硬币结果。节目开始后，参赛者选择一扇门。Monty（知道汽车位置）随后打开一扇门。若汽车被展示，游戏结束；若展示山羊，则参赛者获得换门机会。现在假设参赛者选择门 $1$，随后 Monty 打开门 $2$，展示一只山羊。若参赛者换到门 $3$，成功概率是多少？

  \item 考虑 Monty Hall 问题的如下变体：在某些情形下，Monty 可能\textit{不会}打开门并给参赛者换门机会。具体地，有 $3$ 扇门，其中 $2$ 扇门后有山羊，$1$ 扇门后有汽车。汽车等可能地位于任意门后，Monty 知道汽车位置。令 $0\le p\le 1$。

  参赛者选择一扇门。如果最初选择后有汽车，Monty \textit{会}打开另一扇门展示山羊（在两扇可选门中等可能选择），然后给参赛者是否换到另一扇未打开门的选择。如果最初选择后是山羊，则 Monty 以概率 $p$ \textit{会}打开另一扇门展示山羊，并给参赛者是否换到另一扇未打开门的选择；但以概率 $1-p$，Monty \textit{不会}开门，参赛者必须坚持最初选择。

  参赛者事先决定采用如下策略：初始选择门 $1$。若 Monty 开门并提供换门机会，则换门。
  \begin{enumerate}
    \item 求参赛者得到汽车的无条件概率。并检查在极端情形 $p=0$ 与 $p=1$ 下答案化为什么，简要解释为什么合理。
    \item 现在 Monty 打开门 $2$，展示山羊，因此参赛者换到门 $3$。给定这一信息，求参赛者得到汽车的条件概率。
  \end{enumerate}

  \item 你是 Monty Hall 节目的参赛者。Monty 正在试行新版游戏，规则如下。你可以从三扇门中选一扇。一扇门后有汽车，另一扇后有电脑，另一扇后有山羊（所有排列等可能）。Monty 知道每扇门后的奖品，会打开一扇门（但不是你选的那扇），然后让你选择是否从当前选择换到另一扇未打开的门。

  假设你偏好汽车胜过电脑，电脑胜过山羊，并因此偏好汽车胜过山羊。
  \begin{enumerate}
    \item 仅本小题假设 Monty 总是打开两种替代奖品中你较不喜欢的那个；例如，如果他面对在山羊和电脑之间选择，他会展示山羊。Monty 打开一扇门，展示山羊（仍只对本小题如此）。给定这一信息，你应该换门吗？若换门，得到汽车的成功概率是多少？
    \item 现在假设 Monty 以概率 $p$ 展示你较不喜欢的奖品，以概率 $q=1-p$ 展示你较喜欢的奖品。Monty 打开一扇门，展示电脑。给定这一信息，你应该换门吗（答案可依赖于 $p$）？若换门，得到汽车的成功概率是多少（用 $p$ 表示）？
  \end{enumerate}

  \item Monty Hall 为游戏引入了新变化，把汽车初始位置概率为 $(1/3,1/3,1/3)$ 的假设推广。具体地，有三扇门，其中一扇后有汽车（参赛者想要），另外两扇后有山羊（参赛者不想要）。起初，门 $i$ 后有汽车的概率为 $p_i$，其中 $p_1,p_2,p_3$ 为已知常数，满足 $0<p_1\le p_2\le p_3<1$ 且 $p_1+p_2+p_3=1$。参赛者选择一扇门。随后 Monty 打开一扇门（不是参赛者选择的门），并让参赛者选择是否换到另一扇未打开门。
  \begin{enumerate}
    \item 本小题假设 Monty 事先知道哪扇门后有汽车。他总是打开一扇门展示山羊；若可选择打开哪扇门，则等可能选择。假设本小题中参赛者最初选择门 $3$，随后 Monty 打开门 $2$，展示山羊。给定上述信息，求门 $3$ 后有汽车的条件概率。参赛者是否应该换门？（若是否换门依赖于 $p_i$，请用 $p_i$ 给出完全化简的判别准则。）
    \item 现在改为假设 Monty 事先\textit{不知道}汽车在哪里。他等可能随机选择一扇门打开（不是参赛者选择的门）。（如果他展示汽车，游戏就被毁掉。）再次假设参赛者最初选择门 $3$，随后 Monty 打开门 $2$，展示山羊。给定上述信息，求门 $3$ 后有汽车的条件概率。参赛者是否应该换门？（若是否换门依赖于 $p_i$，请用 $p_i$ 给出完全化简的判别准则。）
    \item 重复 (a)，但参赛者最初选择门 $1$ 而不是门 $3$。
    \item 重复 (b)，但参赛者最初选择门 $1$ 而不是门 $3$。
  \end{enumerate}

  \item Monty Hall 正在试行新版游戏。在这个版本中，不再总是有 $1$ 辆汽车和 $2$ 只山羊；每扇门后的奖品独立生成，每扇门以概率 $p$ 有汽车、以概率 $q=1-p$ 有山羊。参赛者选择一扇门。Monty 知道每扇门内容，随后打开其余两扇门中的一扇。选择打开哪扇门时，若可能他总会展示山羊；若剩余两扇门奖品相同，则等可能随机选择。打开一扇门后，Monty 让参赛者选择是否换到另一扇未打开门。

  参赛者事先决定采用如下策略：先选门 $1$，Monty 开门后换到另一扇未打开门。
  \begin{enumerate}
    \item 求参赛者得到汽车的无条件概率。
    \item 现在 Monty 打开门 $2$，展示山羊。给定这一信息，求参赛者得到汽车的条件概率。
  \end{enumerate}

  \item Monty Hall 试行另一个新版游戏。参赛者从四扇门中选一扇；一扇门后有汽车，一扇后有苹果，一扇后有书，一扇后有山羊，所有 $24$ 种排列等可能。参赛者偏好从低到高为：山羊、苹果、书、汽车。

  Monty 知道每扇门后的奖品，会打开一扇门（不是参赛者最初选择的门），然后让参赛者选择是否换到另一扇未打开门。Monty 以概率 $p$ 展示（在参赛者最初未选中的 $3$ 扇门中）最不受偏好的奖品，以概率 $1-p$ 展示中间偏好的奖品，绝不展示最受偏好的奖品。

  参赛者事先决定采用如下策略：最初选门 $1$。Monty 开门后，在其余两扇未打开门中随机选择一扇换过去（各以概率 $1/2$）。
  \begin{enumerate}
    \item 求参赛者得到汽车的无条件概率。

    提示：以汽车所在位置为条件。
    \item 求 Monty 展示苹果的无条件概率。

    提示：以门 $1$ 后的奖品为条件。
    \item Monty 现在打开一扇门，展示苹果。给定这一信息，求参赛者得到汽车的条件概率。
  \end{enumerate}

  \item 你是 Monty Hall 节目的参赛者。为了让游戏刺激翻倍，Monty 会给你两次换门机会。规则如下：有四扇门，一扇后有汽车，另外三扇后有山羊。起初，汽车所在位置等可能。Monty 知道汽车在哪里；当他可以选择打开哪扇门时，等可能选择。你选择一扇门，为具体起见设为门 $1$。Monty 打开一扇门（不是门 $1$），展示山羊，并让你选择是否换到另一扇门。然后 Monty 再打开另一扇门（不是你当前选择的门），展示另一只山羊。于是现在有两扇开着的山羊门和两扇未打开的门。Monty 再次给你换门机会。

  你事先决定采用以下四种策略之一：留-留、留-换、换-留、换-换。例如“留-换”表示第一次 Monty 给你换门机会时，你保留当前选择；第二次给机会时，你换门。以下各问都要求从游戏开始前的视角求或比较无条件概率。
  \begin{enumerate}
    \item 若采用留-留策略，求赢得汽车的概率。
    \item 若采用留-换策略，求赢得汽车的概率。
    \item 若采用换-留策略，求赢得汽车的概率。
    \item 若采用换-换策略，求赢得汽车的概率。
    \item 四种策略中哪一种最好？
  \end{enumerate}

\subsection*{第一步分析与赌徒破产}

  \item （S）反复掷一枚公平骰子，并记录累计和。令 $p_n$ 为累计和曾经恰好等于 $n$ 的概率（假设骰子会一直掷到累计和最终超过 $n$，但可能从未等于 $n$）。
  \begin{enumerate}
    \item 写出 $p_n$ 的递归方程（以简单方式把 $p_n$ 与较早项 $p_k$ 联系起来）。你的方程应对所有正整数 $n$ 成立，因此请定义 $p_0$ 和 $k<0$ 时的 $p_k$，使小 $n$ 时递归仍成立。
    \item 求 $p_7$。
    \item 直观解释为什么当 $n\to\infty$ 时，$p_n\to 1/3.5=2/7$。
  \end{enumerate}

  \item 进行 $n\ge 1$ 次独立试验，每次试验结果为“成功”或“失败”（不能同时为二者）。令 $p_i$ 为第 $i$ 次试验成功概率，$q_i=1-p_i$，$b_i=q_i-1/2$，其中 $i=1,2,\ldots,n$。令 $A_n$ 为成功次数为偶数的事件。
  \begin{enumerate}
    \item 证明当 $n=2$ 时，$P(A_2)=1/2+2b_1b_2$。
    \item 用归纳法证明
    \[
      P(A_n)=1/2+2^{n-1}b_1b_2\cdots b_n.
    \]
    这个结果在密码学中很有用。它还说明，若抛掷 $n$ 枚硬币，则正面数为偶数的概率为 $1/2$ 当且仅当至少一枚硬币公平。提示：把一些试验合并成一次超级试验。
    \item 直接检查 (b) 的结果在以下简单情形中成立：某个 $i$ 有 $p_i=1/2$；所有 $i$ 都有 $p_i=0$；所有 $i$ 都有 $p_i=1$。
  \end{enumerate}

  \item （S）Calvin 和 Hobbes 进行由一系列游戏组成的比赛，其中 Calvin 每局独立地以概率 $p$ 获胜。比赛采用“赢二”规则：第一个比对手多赢两局的人赢得比赛。分别用两种方法求 Calvin 赢得比赛的概率（用 $p$ 表示）：
  \begin{enumerate}
    \item 以条件化和全概率公式求解；
    \item 把问题解释为赌徒破产问题。
  \end{enumerate}

  \item （S）一名赌徒反复玩一个游戏，每轮以概率 $1/3$ 赢 $1$ 美元，以概率 $2/3$ 输 $1$ 美元。他的策略是“领先 $2$ 美元时退出”。假设他初始有一百万美元。证明他曾经领先 $2$ 美元的概率小于 $1/4$。

  \item 如赌徒破产问题中，两名赌徒 A 与 B 进行一系列赌博，直到其中一人破产。令 A 初始有 $i$ 美元，B 初始有 $N-i$ 美元，A 赢得每局的概率为 $p$，其中 $0<p<1/2$。每局赌注为 $1/k$ 美元，其中 $k$ 是正整数；例如，$k=1$ 是原始赌徒破产问题，$k=20$ 表示他们赌五美分。求 A 赢得游戏的概率，并判断当 $k\to\infty$ 时会发生什么。

  \item 圆周上有 $100$ 个等距点。$99$ 个点上有羊，$1$ 个点上有狼。每个时刻，狼随机顺时针或逆时针移动一个点。若该点有羊，狼就吃掉它。羊不移动。初始时位于狼对面的那只羊最后存活的概率是多少？

  \item 一个不死的醉汉在整数线上随机游走。他从原点开始，每一步以概率 $p$ 向右移动 $1$ 个单位，以概率 $q=1-p$ 向左移动 $1$ 个单位，且各步相互独立。令 $S_n$ 为他第 $n$ 步后的位置。
  \begin{enumerate}
    \item 令 $p_k$ 为醉汉曾经到达值 $k$ 的概率，对所有 $k\ge 0$。写出 $p_k$ 的差分方程（本小题无需求解）。
    \item 求完全化简的 $p_k$；务必考虑三种情形：$p<1/2,\ p=1/2,\ p>1/2$。可以使用如下事实：若 $A_1,A_2,\ldots$ 满足 $A_j\subseteq A_{j+1}$ 对所有 $j$ 成立，则 $P(A_n)\to P(\cup_{j=1}^\infty A_j)$（概率的连续性）。
  \end{enumerate}

\subsection*{Simpson 悖论}

  \item （S）
  \begin{enumerate}
    \item 是否可能存在事件 $A,B,C$，使得 $P(A|C)<P(B|C)$ 且 $P(A|C^c)<P(B|C^c)$，但 $P(A)>P(B)$？也就是说，给定 $C$ 为真时 $A$ 比 $B$ 不可能，给定 $C$ 为假时 $A$ 也比 $B$ 不可能，但在没有关于 $C$ 的信息时 $A$ 比 $B$ 更可能。证明这是可能的（用简短证明），或给出一个反例（用故事解释 $A,B,C$）。
    \item 若 (a) 的情形可能，它是 Simpson 悖论的特例、等价于 Simpson 悖论，还是二者都不是？若不可能，直观解释为什么虽然 Simpson 悖论可能发生，这里却不可能。
  \end{enumerate}

  \item （S）考虑《辛普森一家》某集中的对话：
  \begin{quote}
  Lisa：爸爸，我觉得他是象牙贩子！他的靴子是象牙的，帽子是象牙的，我很确定那张支票也是象牙的。

  Homer：Lisa，一个有很多象牙的人，比象牙库存少的人更不可能伤害 Stampy。
  \end{quote}
  这里 Homer 和 Lisa 正在争论：如果他们把大象 Stampy 卖给名叫 Blackheart 的人，他是否可能伤害 Stampy。他们显然对于如何利用关于 Blackheart 的观察，来了解给定证据后 Blackheart 会伤害 Stampy 的概率，有不同看法。
  \begin{enumerate}
    \item 为这里的各种相关事件定义清楚记号。
    \item 用 (a) 中的记号，把 Lisa 和 Homer 的论点（Lisa 的论点部分是隐含的）表示为条件概率陈述。
    \item 假设“某人拥有大量某商品时，获得更多该商品的欲望较低”是真的。解释 Homer 的推理为什么仍然有问题：关于 Blackheart 的证据为什么不能说明他伤害 Stampy 的可能性更低。
  \end{enumerate}

  \item
  \begin{enumerate}
    \item 有两个深红色罐子（标记为 $C_1,C_2$）和两个淡紫色罐子（标记为 $M_1,M_2$）。每个罐子中都有绿色软糖熊和红色软糖熊的混合物。举例说明：可能出现 $C_1$ 中绿色软糖熊比例远高于 $M_1$，且 $C_2$ 中绿色软糖熊比例远高于 $M_2$，但若把 $C_1,C_2$ 合并成一个新罐子、也把 $M_1,M_2$ 合并成一个新罐子，则 $C_1,C_2$ 的合并罐绿色比例低于 $M_1,M_2$ 的合并罐。
    \item 解释 (a) 如何与 Simpson 悖论相关：既给出直观说明，也在 Simpson 悖论的陈述中明确地定义事件 $A,B,C$。
  \end{enumerate}

  \item 如本章所述，Simpson 悖论说可能存在事件 $A,B,C$，使得
  \[
    P(A|B,C)<P(A|B^c,C),\qquad
    P(A|B,C^c)<P(A|B^c,C^c),
  \]
  但
  \[
    P(A|B)>P(A|B^c).
  \]
  \begin{enumerate}
    \item 若 $A$ 与 $B$ 独立，Simpson 悖论能否发生？若能，给出具体例子（包括数字和解释）；若不能，证明不可能。
    \item 若 $A$ 与 $C$ 独立，Simpson 悖论能否发生？若能，给出具体例子；若不能，证明不可能。
    \item 若 $B$ 与 $C$ 独立，Simpson 悖论能否发生？若能，给出具体例子；若不能，证明不可能。
  \end{enumerate}

  \item （S）Andrew Gelman 的书《Red State, Blue State, Rich State, Poor State》讨论了如下选举现象：在美国任意州内，富裕选民比贫穷选民更可能投票给共和党；然而，更富裕的州却倾向于支持民主党候选人。
  \begin{enumerate}
    \item 为简单起见，假设只有两个州（称为红州和蓝州），每州 $100$ 人，且每个人要么富裕要么贫穷，要么是民主党人要么是共和党人。构造与上述现象一致的数字，为每个州给出一个 $2\times 2$ 表（列出每州中富裕民主党人等人数）。也就是说，在每个州内，富裕选民比贫穷选民更可能投给共和党，但富裕人口比例较高的州民主党比例更高。
    \item 在 (a) 的设定中（不一定使用你编造的数字），令 $D$ 为随机选一人是民主党人的事件（所有 $200$ 人等可能），令 $B$ 为此人住在蓝州的事件。假设 $10$ 人从蓝州搬到红州。用 $P_{\rm old}$ 与 $P_{\rm new}$ 分别表示搬迁前后的概率。假设人们不改变党派，因此 $P_{\rm new}(D)=P_{\rm old}(D)$。是否可能同时有
    \[
      P_{\rm new}(D|B)>P_{\rm old}(D|B),\qquad
      P_{\rm new}(D|B^c)>P_{\rm old}(D|B^c)?
    \]
    若可能，解释这为什么不与全概率公式矛盾；若不可能，证明不可能。
  \end{enumerate}

\subsection*{综合练习}

  \item 一名患者接受 conditionitis 疾病的血液检测。令 $p$ 为患者患 conditionitis 的先验概率。血样被送往两个实验室之一分析，即实验室 A 或实验室 B。使用哪个实验室的选择随机进行，与患者患病状态独立，且两个实验室概率各为 $1/2$。

  对实验室 A，给定某人有病时检测为阳性的概率为 $a_1$，给定某人无病时检测为阴性的概率为 $a_2$。实验室 B 的对应概率为 $b_1,b_2$。
  \begin{enumerate}
    \item 给定检测为阳性，求患者患病的概率。
    \item 给定检测为阳性，求患者血样由实验室 A 分析的概率。
  \end{enumerate}

  \item Fred 决定进行一系列 $n$ 次检测，以诊断他是否患有某疾病（任何单个检测都不完全可靠，所以他希望通过多次检测降低不确定性）。令 $D$ 为他患病的事件，$p=P(D)$ 为先验概率，$q=1-p$。令 $T_j$ 为第 $j$ 次检测为阳性的事件。
  \begin{enumerate}
    \item 本小题假设给定 Fred 的患病状态后，检测结果条件独立。令 $a=P(T_j|D)$，$b=P(T_j|D^c)$，其中 $a,b$ 不依赖于 $j$。给定 $n$ 次检测全为阳性，求 Fred 患病的后验概率。
    \item 假设 Fred 的 $n$ 次检测全为阳性。然而，有些人带有某基因，使他们\textit{总是}检测为阳性。令 $G$ 为 Fred 带有该基因的事件。假设 $P(G)=1/2$，且 $D$ 与 $G$ 独立。若 Fred 没有该基因，则给定患病状态后，检测结果条件独立。令 $a_0=P(T_j|D,G^c)$，$b_0=P(T_j|D^c,G^c)$，其中 $a_0,b_0$ 不依赖于 $j$。给定 $n$ 次检测全为阳性，求 Fred 患病的后验概率。
  \end{enumerate}

  \item 某种遗传病可以由母亲传给孩子。给定母亲患病，她的每个孩子独立地以概率 $1/2$ 患病；给定母亲不患病，她的孩子也不会患病。某位母亲患病概率为 $1/3$，她有两个孩子。
  \begin{enumerate}
    \item 求两个孩子都不患病的概率。
    \item 年长孩子是否患病，与年幼孩子是否患病独立吗？解释。
    \item 发现年长孩子不患病。一周后，又发现年幼孩子也不患病。给定这些信息，求母亲患病的概率。
  \end{enumerate}

  \item 同时抛三枚公平硬币。解释下面论证的问题：“三枚硬币全同面的概率是 $50\%$，因为显然总能找到两枚匹配的硬币，然后另一枚硬币有 $50\%$ 的概率与这两枚匹配。”

  \item 一个坛中有红、绿、蓝三色球。令 $r,g,b$ 分别为红、绿、蓝球比例，且 $r+g+b=1$。
  \begin{enumerate}
    \item 有放回地随机抽球。求第一次抽到绿球早于第一次抽到蓝球的概率。

    提示：解释这与“给定一次抽取为绿球或蓝球，求它为绿球的概率”之间的关系。
    \item 无放回地随机抽球。求第一次抽到绿球早于第一次抽到蓝球的概率。答案与 (a) 相同还是不同？

    提示：想象所有球按将被抽出的顺序排成一列。注意红球在队列中的位置无关紧要。
    \item 将 (a) 的结果推广到如下设定：进行独立试验，每次结果恰好归入类型 $1,2,\ldots,n$ 之一，概率分别为 $p_1,p_2,\ldots,p_n$。求当 $i\ne j$ 时，第一次出现类型 $i$ 早于第一次出现类型 $j$ 的概率。
  \end{enumerate}

  \item Marilyn vos Savant 在《Parade》专栏中被问到如下问题：
  \begin{quote}
  你在一个聚会上，另有 $199$ 位客人。此时劫匪闯入，宣布要抢劫其中一人。他们把 $199$ 张空白纸条和 $1$ 张写着“you lose”的纸条放进帽子，每位客人必须抽一张，抽到“you lose”的人会被抢。劫匪允许你选择第一个抽、最后一个抽，或在中间任意时刻抽。你会什么时候抽？
  \end{quote}
  抽取是不放回的，且 (a) 中各纸条均匀随机。
  \begin{enumerate}
    \item 判断为了最大化不被抢的概率，最优是第一个抽、最后一个抽、中间某时刻抽，还是无所谓。给出清楚、简洁且有说服力的解释。
    \item 更一般地，假设有 $1$ 张“you lose”纸条，权重为 $v$；另有 $n$ 张空白纸条，每张权重为 $w$。每一阶段抽取按权重成比例进行，即抽到某张纸条的概率等于其权重除以所有剩余纸条权重之和。判断第一个抽和第二个抽哪个更好（或是否无所谓）；其中 $v>0,w>0,n\ge 1$ 为已知常数。
  \end{enumerate}

  \item 反复掷一枚公平骰子，直到累计和至少为 $100$（此时停止）。求最终累计和最可能的值，即累计和第一次至少为 $100$ 时的值。

  提示：考虑最后一次掷骰前累计和可能是多少。

  \item Homer 有一盒甜甜圈，目前恰有 $c$ 个巧克力味、$g$ 个糖霜味、$j$ 个果冻味。Homer 逐个吃甜甜圈，每次从剩余甜甜圈中均匀随机选择。
  \begin{enumerate}
    \item 求盒中最后剩下的甜甜圈是巧克力味的概率。
    \item 求以下事件的概率：糖霜味是 Homer 第一个吃完的甜甜圈类型，然后果冻味是第二个吃完的类型。

    提示：考虑最后剩下的甜甜圈，以及糖霜或果冻中最后剩下的那个甜甜圈。
  \end{enumerate}

  \item 令 $D$ 为某人患某疾病的事件，$C$ 为此人暴露于某种物质的事件（例如 $D$ 可以对应肺癌，$C$ 可以对应吸烟）。我们关心暴露是否与患病相关，以及若相关，它们如何相关。流行病学中广泛使用\textbf{赔率比}来度量疾病与暴露之间的关联，定义为
  \[
    \operatorname{OR}=\frac{\operatorname{odds}(D|C)}
                           {\operatorname{odds}(D|C^c)},
  \]
  其中条件赔率类似于无条件赔率：
  \[
    \operatorname{odds}(A|B)=\frac{P(A|B)}{P(A^c|B)}.
  \]
  另一种广泛使用的度量是暴露者相对于未暴露者的\textbf{相对风险}
  \[
    \operatorname{RR}=\frac{P(D|C)}{P(D|C^c)}.
  \]
  相对风险尤其容易解释，例如 $\operatorname{RR}=2$ 表示暴露于该物质的人患病概率是未暴露者的两倍（但这未必表示该物质导致患病概率上升，赔率比也未必有因果解释）。
  \begin{enumerate}
    \item 证明：若疾病很罕见，无论对暴露者还是未暴露者都是如此，那么相对风险近似等于赔率比。
    \item 令 $p_{ij}$（$i=0,1$ 且 $j=0,1$）为下列 $2\times 2$ 表中的概率：
    \[
    \begin{array}{c|cc}
      & D & D^c\\
      \hline
      C & p_{11} & p_{10}\\
      C^c & p_{01} & p_{00}
    \end{array}
    \]
    例如，$p_{10}=P(C,D^c)$。证明赔率比可表示为交叉乘积比：
    \[
      \operatorname{OR}=\frac{p_{11}p_{00}}{p_{10}p_{01}}.
    \]
    \item 证明赔率比具有漂亮的对称性：交换 $C$ 与 $D$ 的角色不会改变其值：
    \[
      \operatorname{OR}=\frac{\operatorname{odds}(C|D)}
                              {\operatorname{odds}(C|D^c)}.
    \]
    这个性质是赔率比被广泛使用的主要原因之一，因为它使得赔率比可以在相对风险难以很好估计的各种问题中被估计。
  \end{enumerate}

  \item 一位研究者想通过调查估计某总体中使用过非法药物的人所占比例。由于担心很多人在被直接问“你是否使用过非法药物？”时会撒谎，研究者使用一种称为\textbf{随机化回答}的方法。帽子中装有纸条，每张纸条写着“我使用过非法药物”或“我没有使用过非法药物”。令 $p$ 为写着“我使用过非法药物”的纸条比例（$p$ 由研究者预先选择）。

  每位参与者从帽子中随机抽一张纸条，并如实回答该纸条上的陈述是否为真（“是”或“否”）。纸条随后放回帽子。研究者不知道参与者抽到哪种纸条。令 $y$ 为参与者回答“是”的概率，$d$ 为参与者使用过非法药物的概率。
  \begin{enumerate}
    \item 用 $d,p$ 表示 $y$。
    \item 研究者在设计调查时可能作出的最坏 $p$ 值选择是什么？解释。
    \item 现在考虑另一种系统。假设比例 $p$ 的纸条写着“我使用过非法药物”，但其余 $1-p$ 的纸条现在写着“我出生在冬季”，而不是“我没有使用过非法药物”。假设 $1/4$ 的人出生在冬季，且出生季节与是否使用过非法药物独立。用 $y,p$ 表示 $d$。
  \end{enumerate}

  \item 在 Tom Stoppard 的戏剧《Rosencrantz and Guildenstern Are Dead》开头，Guildenstern 在抛硬币，Rosencrantz 对每次结果下注。硬币一次又一次落到正面，于是出现如下评论：
  \begin{quote}
  Guildenstern：如果没有别的，弱者也许会重新审视他的信念，至少会重新审视概率法则。
  \end{quote}
  硬币连续 $92$ 次为正面。
  \begin{enumerate}
    \item Fred 和朋友正在看这出戏。看到上述事件后，他们有如下对话：
    \begin{quote}
    Fred：如果硬币公平，这个结果极不可能。他们一定在用有问题的硬币（也许是双正面硬币），或者实验一定以某种方式被操纵了（也许用了磁铁）。

    Fred 的朋友：长度为 $92$ 的全正面序列确实极不可能；公平硬币下概率为 $1/2^{92}\approx 2\times 10^{-28}$。但任意其他具体的长度为 $92$ 的 H/T 序列概率完全相同！全正面看起来极不可能的原因，只是可能结果数随抛掷次数指数增长，所以任何结果都会显得极不可能。你甚至可以在不看实验结果的情况下提出同样论证，这说明你实际上没有反对硬币公平的证据。
    \end{quote}
    讨论这些评论，帮助 Fred 和他的朋友解决争论。
    \item 假设只有两种可能：硬币全都公平（且抛掷公平），或使用双正面硬币（此时正面概率为 $1$）。令 $p$ 为硬币公平的先验概率。给定 $92$ 次中 $92$ 次都为正面，求硬币公平的后验概率。
    \item 延续 (b)，哪些 $p$ 值使硬币公平的后验概率大于 $0.5$？哪些 $p$ 值使它小于 $0.05$？
  \end{enumerate}

  \item 有 $n$ 种玩具，你逐个收集。每次购买一个玩具时，其类型等可能随机确定。令 $p_{ij}$ 为购买第 $i$ 个玩具后，你的收藏中恰有 $j$ 种玩具的概率，其中 $i\ge 1,\ 0\le j\le n$。（这是优惠券收集问题的背景；我们会在例 4.3.12 中研究这个著名问题。）
  \begin{enumerate}
    \item 对 $i\ge 2$ 且 $1\le j\le n$，找出用 $p_{i-1,j}$ 与 $p_{i-1,j-1}$ 表示 $p_{ij}$ 的递归方程。
    \item 说明如何用 (a) 中的递归计算 $p_{ij}$。
  \end{enumerate}

  \item A/B 测试是许多公司用来了解顾客对不同处理方式反应的随机化实验。例如，公司可能想了解用户对网站新功能的反应，或比较两个广告。顾名思义，研究两种处理方式 A 与 B。用户逐个到来，每位用户被随机分配到两种处理之一。每位用户的试验结果被分类为“成功”（例如完成购买）或“失败”。第 $n$ 个用户接受处理 A 的概率允许依赖于先前用户的结果；这种设定称为\textbf{双臂老虎机}。

  人们研究了许多随机分配处理方式的算法。下面是一种特别简单（但反复无常）的算法，称为“跟随胜者”过程：
  \begin{enumerate}
    \item[(i)] 第一个用户等可能随机分配到处理 A 或处理 B。
    \item[(ii)] 若第 $n$ 个用户的试验成功，则第 $n+1$ 个用户继续使用同一处理；否则，切换到另一个处理。
  \end{enumerate}
  令 $a$ 为处理 A 的成功概率，$b$ 为处理 B 的成功概率。假设 $a\ne b$，但 $a,b$ 未知（这正是需要测试的原因）。令 $p_n$ 为第 $n$ 次试验成功的概率，$a_n$ 为第 $n$ 次试验分配到处理 A 的概率（使用上述算法）。
  \begin{enumerate}
    \item 证明
    \[
      p_n=(a-b)a_n+b,\qquad
      a_{n+1}=(a+b-1)a_n+1-b.
    \]
    \item 利用 (a) 的结果证明 $p_{n+1}$ 满足递归方程
    \[
      p_{n+1}=(a+b-1)p_n+a+b-2ab.
    \]
    \item 假设极限存在，利用 (b) 求该算法长期成功概率 $\lim_{n\to\infty}p_n$。
  \end{enumerate}
\end{enumerate}

\end{document}
